\documentclass[11pt]{article}
\usepackage{amsmath,amssymb,amsthm}
\setlength{\textwidth}{6.3in}
\setlength{\oddsidemargin}{0.1in}
\setlength{\evensidemargin}{0.1in}
\setlength{\textheight}{8.6in}
\setlength{\topmargin}{-0.3in}
\setlength{\headheight}{14pt}
\setlength{\headsep}{22pt}
\setlength{\parindent}{1.2em}
\setlength{\parskip}{3pt}
\newcommand{\R}{\mathbb R}
\newcommand{\N}{\mathbb N}
\newcommand{\Nzero}{\mathbb N_0}
\newcommand{\gra}{\operatorname{gra}}
\newcommand{\dom}{\operatorname{dom}}
\newcommand{\sgn}{\operatorname{sgn}}
\newcommand{\pair}[2]{\langle #1,#2\rangle}
\newtheorem{theorem}{Theorem}
\newtheorem{proposition}{Proposition}
\newtheorem{lemma}{Lemma}
\newtheorem{corollary}{Corollary}
\renewcommand{\thecorollary}{C.\arabic{corollary}}
\newcommand{\dist}{\operatorname{dist}}
\pagestyle{plain}
\begin{document}

\begin{center}
{\Large\bfseries Nonmaximal sums of maximally monotone operators}\\[10pt]
\ifdefined\PrivateAuthorBlock\PrivateAuthorBlock\\[10pt]\fi
September 7, 2026\\[4pt]
\ifdefined\expandedversion Extended manuscript\else Main manuscript\fi
\end{center}


\begin{abstract}
\ifdefined\expandedversion
We construct two maximally monotone operators on a real Banach space
whose domains satisfy Rockafellar's interior-domain condition and whose
sum is not maximally monotone. In the determined construction, the second
operator is bounded, linear,
positive, rank one, and defined on the whole space. We first construct
the operators on $c_0$ and then realize the same conclusion on standard
$\ell^1$ with its full dual $\ell^\infty$. The construction uses countably
many copies of a classical triangular operator and constrains their sums
by a Lipschitz curve. Every parameter in this curve has a finitely
supported tail, and its limiting tail belongs to $\ell^2\setminus\ell^1$.
We prove maximality by reducing every point of the monotone polar to a
scalar inequality and excluding the missing parameter interval through
this failure of summability. The first operator omits zero and satisfies
a uniform lower bound on the pairing divided by the primal norm.
The positive rank-one perturbation makes the pairing strictly positive
on every point of the sum graph, so that $(0,0)$ gives a proper monotone
extension. For the same first factor, every positive multiple of the
rank-one perturbation gives a nonmaximal sum, including arbitrarily small
perturbations. We also compute a positive Fitzpatrick split value at the origin.
An abstract endpoint criterion characterizes maximality within the same
construction and expresses the entire monotone polar using actual dual
labels. It identifies the exact condition that the nonsummable tail verifies.
We also construct a countable monotone seed and prove that it admits a
maximal extension preserving a lower pairing bound. Every selected
extension differs from the original first factor and gives a nonmaximal
sum with a full-domain norm subdifferential, on both spaces. This second
family is nonconstructive, with all extension premises verified explicitly.
The extended version also treats normal cones to bounded convex sets,
derives consequences for type (FPV) and regularity, and studies Lipschitz
extensions and convex combinations of graph points.
\else
We construct a maximally monotone operator on $c_0$ whose sum with a
bounded positive rank-one operator is not maximally monotone, although
the second operator has full domain and Rockafellar's interior-domain
condition holds. The construction couples countably many triangular
operators through a Lipschitz curve of block sums. Each actual parameter
has a finitely supported tail; the common limiting tail belongs to
$\ell^2\setminus\ell^1$. We compute the entire monotone polar and prove
maximality by showing that a point over the missing parameter interval
would require a nonsummable continuous-dual label. An exact pairing
identity then supplies a proper monotone extension of the sum. We also
prove an abstract criterion that identifies the precise failure of
endpoint realization used in the construction, and derive the finite
radial bound and the nonconvexity of the first operator's domain closure.
\fi
\end{abstract}

\noindent\textbf{Keywords:} maximally monotone operator; sum theorem;
nonreflexive Banach space; monotone polar; rank-one operator.

\section{Introduction}
\label{sec:introduction}
Let $X$ be a real Banach space, let $X^*$ be its continuous dual, and let
$A,B:X\rightrightarrows X^*$ be maximally monotone operators. Their pointwise
sum is defined by $(A+B)x=\{a+b:a\in Ax,\ b\in Bx\}$. Rockafellar's sum
theorem \cite{rockafellar1970} states that $A+B$ is maximally monotone when
$X$ is reflexive and
\begin{equation}
\dom A\cap\operatorname{int}(\dom B)\ne\varnothing.
\label{eq:original-cq}
\end{equation}
The unrestricted sum question asks whether Eq.~\eqref{eq:original-cq}
alone is sufficient on an arbitrary real Banach space. The interior here
is taken in the norm topology of $X$, and maximality is taken in the full
pair $X\times X^*$.
The broader theory of maximality and convex representation on Banach
spaces is developed in Simons' monograph \cite{simons2008}.

Several positive results extend the theorem by imposing structure on one
of the factors. Yao \cite{yao2010} proves maximality when the first factor
is of type~(FPV) and the second factor has full domain. Borwein and Yao
\cite{borweinyao2012} prove the sum theorem under
Eq.~\eqref{eq:original-cq} when the first factor is a linear relation.
The factor order in the latter qualification matters: placing a
full-domain linear operator first would require interior points in the
domain of the other factor. These results motivate the question of
whether an elementary full-domain second factor can already produce a
nonmaximal sum with a nonlinear first factor.
Normal-cone partners under type~(FPV) hypotheses are treated by Voisei
\cite{voisei2010}. At the universal level, Bo\c{t}, Bueno, and Simons
\cite{botbuenosimons2016} show that restricting the sum question to pairs
with bounded common domain does not weaken it. This is a statement about
all admissible pairs, not a localization theorem for one fixed factor.

Convex representations provide a second way to organize these questions.
Fitzpatrick \cite{fitzpatrick1988} associates a convex function with a
monotone graph. Burachik and Svaiter \cite{burachiksvaiter2002} study the
associated family of convex representatives through enlargements, and
Penot \cite{penot2004} develops representation methods for compositions
and sums in the reflexive setting. For nonmaximal graphs,
representability and maximality must still be distinguished;
Mart\'{i}nez-Legaz and Svaiter \cite{martinezsvaiter2005} study this
distinction, including unique maximal extensions in finite dimensions.
Bueno, Mart\'{i}nez-Legaz,
and Svaiter \cite{buenomartinezsvaiter2014} compare representable and
monotone-polar closures. Our endpoint criterion concerns graphs containing
every point in specified affine fibres and computes their polar directly; it does not
replace these representation theories by an unrestricted normal form.
Marques Alves and Svaiter \cite{marquessvaiter2012} prove a sum theorem
under a qualification on domains of Fitzpatrick representations.
Such a qualification must be checked on those representations; it cannot
be replaced without proof by Eq.~\eqref{eq:original-cq}.

Nonreflexive constructions provide a concrete setting for the sum question.
Gossez's work \cite{gossez1971} develops monotone-operator theory beyond
reflexive spaces. The dense-type class has additional structure:
Marques Alves and Svaiter \cite{marquessvaiter2010} identify types~(D)
and~(NI), and Bauschke, Borwein, Wang, and Yao \cite{bauschkefp2012}
prove that Fitzpatrick--Phelps type~(FP) implies dense type.
Type~(FP) here is not type~(FPV). These results distinguish important
subclasses; none of these type assumptions is imposed in the question above.
Bueno and Svaiter \cite{buenosvaiter2011} construct a non-type~(D)
maximally monotone operator on $c_0$. Bauschke, Borwein, Wang, and Yao
\cite{bauschkeetal2011} develop a family of triangular operators and
further nonlinear examples. In particular, Example~4.1 of their
August~6, 2011 author version contains
the single-block triangular map used below. These earlier ingredients
give exact pairing identities. To obtain the sum behavior considered
here, we impose an additional constraint on the sums over countably many
blocks and verify maximality for every point of the full monotone polar.

We define this constraint through a $1$-Lipschitz map
$F:J\to\ell^1(\Nzero)\subset\ell^2(\Nzero)$, where
$J=(-\infty,-1)\cup(1,\infty)$. Its positive-index coordinates have finite
support for each actual parameter. At both deleted endpoints they tend,
coordinatewise and in $\ell^2$, to
$(1/(4n))_{n\ge1}$, which is not in $\ell^1$.
The triangular pairing identity first proves monotonicity. The complete
affine fibres of the graph then reduce the condition for an arbitrary
polar point to
\begin{equation}
\|F(t)-m_p\|_2^2+\nu^2\le(t-\tau)^2
\qquad(t\in J),
\label{eq:polar-intro}
\end{equation}
where $m_p\in\ell^1(\Nzero)$ is determined by its dual coordinate and
$\tau,\nu\in\R$. If $|\tau|>1$, the actual test $t=\tau$ returns a graph
point. If $|\tau|\le1$, finite-coordinate limits at the two endpoints
force the positive-index coordinates of $m_p$ to equal $1/(4n)$.
Their nonsummability excludes this second case and proves maximality.

The resulting first factor $A$ has nonempty domain at positive distance
from zero. Its construction also gives
$\langle x,a\rangle>-\sigma$ on $\gra A$, with
$\sigma=\pi^2/96<1/8$. We choose
$Px=\langle x,g\rangle g$ for a specified nonzero $g\in(c_0)^*$.
This factor is also the subdifferential of the continuous convex quadratic
$x\mapsto\langle x,g\rangle^2/2$, consistent with Rockafellar's
subdifferential maximality theorem \cite{rockafellar1970subdiff}.
We give a direct maximality proof for $P$ below.
The same pairing identity gives
$\langle x,a+Px\rangle>1-\sigma$ on every sum row.
Thus $(0,0)$ is monotonically related to the entire sum and is not in
its graph. The full domain of $P$ makes
Eq.~\eqref{eq:original-cq} immediate from nonemptiness of $\dom A$.

The construction on $c_0$ already addresses the unrestricted question.
\ifdefined\expandedversion
We also transfer it to standard $\ell^1$ through an explicitly specified
surjection $Q:\ell^1\to c_0$. The pullback $T=Q^*AQ$ is proved maximally
monotone in $\ell^1\times\ell^\infty$: tests along $\ker Q$ force every
polar dual coordinate to factor through $Q$, and the remaining inequality
is exactly the one for $A$. This argument uses a controlled preimage for
each point, not a bounded linear right inverse.
These are two realizations of one construction, rather than two independent
mechanisms.

A second family uses a genuinely unspecified maximal extension.
We retain countably many tests in kernel directions approaching the two endpoints
and add a compatible row that is not in the original graph. The resulting
seed excludes every dual label over a fixed primal slab and satisfies a
uniform lower pairing bound. A radial argument proves that an extension
maximal under this bound is globally maximally monotone. Every selected
factor differs from the determined one, while a full-domain norm
subdifferential makes its sum nonmaximal. Theorem~\ref{thm:nonconstructive}
proves the seed conditions and the complete certificates on both spaces;
it does not assume an unconstructed seed or a safe maximal extension.

The same construction also yields nonmaximality for every positive
multiple of the rank-one factor. For strengths below $\sigma$, the origin
is no longer a polar point; a finite-support query on $c_0$ supplies the required
proper monotone extension instead. We prove this parameter statement
on both spaces in Section~\ref{sec:parameters}.

Section~\ref{sec:endpoint-criterion} states and proves the abstract
endpoint criterion behind this construction. It distinguishes a Lipschitz
extension of the curve in the Hilbert space from the existence of elements
of the continuous dual satisfying its constraints.
After the definitions in Section~\ref{sec:definitions}, we construct the
counterexample on $c_0$ in Section~\ref{sec:c0} and transfer it to standard
$\ell^1$ in Section~\ref{sec:l1}.
Section~\ref{sec:nonconstructive} proves the nonconstructive family,
and Section~\ref{sec:structure} explains the representation and the
common mechanism, including convergent primals with unbounded dual labels.
Section~\ref{sec:domain-comparison} records the domain geometry and its
relation to a published almost-convexity statement.
Appendix~\ref{sec:split} proves
$\Delta_{A,P}(0,0)=\Delta_{T,B}(0,0)=\sigma$, each uniquely minimized
at the zero dual split. These identities refine the counterexamples;
they are not needed for the direct nonmaximality proofs.
\ifdefined\expandedversion
The remaining appendices derive the norm-subdifferential and bounded
normal-cone partners and distinguish them from parameter changes,
transported examples, and logical consequences of the same first factor.
\fi
\else
The initial manuscript contains the complete coordinate construction
and an abstract criterion for its maximality. The criterion explains
why completion of the block-sum comparison space need not supply an
element of the continuous dual. Its assumptions are verified by the
coordinate proof, so no representation theorem for arbitrary monotone
operators is required.

After the definitions in Section~\ref{sec:definitions},
Theorem~\ref{thm:C} and its proof give the full counterexample.
Section~\ref{sec:endpoint-criterion} characterizes its endpoint
mechanism, Section~\ref{sec:domain-comparison} treats the domain closure
and the relevant published statements, and Section~\ref{sec:structure}
explains the roles of the affine fibres, the nonsummable limit and the
second factor. The standard-$\ell^1$ transfer, nonconstructive extensions
of the same seed mechanism, parameter variations and further corollaries
are retained in the extended manuscript.
\fi

\section{Definitions}
\label{sec:definitions}
All scalars are real. Put $\N=\{1,2,\ldots\}$, $\Nzero=\{0,1,\ldots\}$,
and $I=\Nzero\times\N$. The space $Y=c_0(I)$ is global: for each
$\epsilon>0$, only finitely many coordinates have modulus at least
$\epsilon$. It has norm $\|\cdot\|_\infty$ and full dual $\ell^1(I)$,
with $\pair{x}{a}=\sum_{b,j}x_{b,j}a_{b,j}$.
The same pairing is used when the first argument lies in $\ell^\infty(I)$.
The unit coordinate vectors are $e_{b,j}$; the ambient space follows their use.

For a graph $G\subset X\times X^*$, define
\begin{equation}
G^\mu=\{(z,p)\in X\times X^*:
 \pair{z-x}{p-a}\ge0\quad\forall(x,a)\in G\}.
\tag{D1}\label{eq:D1}
\end{equation}
The graph is monotone when every two graph points have a nonnegative bracket,
and maximally monotone when it has no proper monotone extension in this same
full dual pair. For a monotone $G$, $G=G^\mu$ is equivalent to maximality.
Domain and range refer to actual graph points. Every inequality quantified
over the graph includes every value at a multivalued primal.

\ifdefined\expandedversion
The convex function $F_A$ below is the Fitzpatrick function
\cite{fitzpatrick1988}; subtracting the duality pairing fixes the gap
sign convention used in this paper.
For a maximally monotone $A$, the sign conventions are
\begin{equation}
\begin{aligned}
\Phi_A(d,p)&=F_A(d,p)-\pair{d}{p}
 =\sup_{(y,a)\in\gra A}\pair{d-y}{a-p},\\
\Delta_{A,B}(x,x^*)&=\inf_{u^*\in X^*}
 [\Phi_A(x,x^*-u^*)+\Phi_B(x,u^*)].
\end{aligned}
\tag{D2}\label{eq:D2}
\end{equation}
Equivalently, $\Delta_{A,B}(x,x^*)=(F_A\square_2 F_B)(x,x^*)
-\pair{x}{x^*}$, where $\square_2$ denotes partial infimal convolution
in the dual variable: $(F_A\square_2 F_B)(x,x^*)$
is the infimum of $F_A(x,x^*-u^*)+F_B(x,u^*)$ over $u^*\in X^*$.
This is not an identification with the Fitzpatrick function of the sum.
\fi
For the present graphs, which have no zero primal, define
\begin{equation}
\mathcal V(G)=\sup_{(x,a)\in G}
 \frac{\max\{0,-\pair{x}{a}\}}{\|x\|}.
\tag{D3}\label{eq:D3}
\end{equation}
At $(0,0)$ this is the nonnegative slope used by Verona and Verona
\cite{verona2008}. The finite upper bounds proved below are properties
of the constructed graphs, not an assumed general bound from that source.

We index coordinates by blocks. The map $m$ records the sum in each block,
and $E$ applies a triangular operator within that block. Its single-block
formula is the $\alpha=(1,1,\ldots)$ case of
\cite[author version, Example~4.1]{bauschkeetal2011}; the additional constraint below
couples the block sums through $F$.
For $a\in\ell^1(I)$, define the bounded maps
\begin{equation}
(ma)_b=\sum_{j\ge1}a_{b,j},\qquad
(Ea)_{b,j}=a_{b,j}+2\sum_{k>j}a_{b,k}.
\tag{D4}\label{eq:D4}
\end{equation}
Here $m:\ell^1(I)\to\ell^1(\Nzero)\subset\ell^2(\Nzero)$ and
$E:\ell^1(I)\to Y$. Put
\begin{equation}
\begin{aligned}
g&=e_{0,1}-e_{0,2}\in Y^*,&
h&=Eg=-e_{0,1}-e_{0,2}\in Y,\\
r(a)&=\pair{h}{a}=-a_{0,1}-a_{0,2},&
La&=-Ea+r(a)h.
\end{aligned}
\tag{D5}\label{eq:D5}
\end{equation}

For $s>0$ and $n\ge1$, write
\begin{equation}
\omega_s(n)=\frac{\max\{1-sn^{1/4},0\}}{4n},\quad
W(s)=\sum_{n\ge1}\omega_s(n)^2,\quad
\sigma=\sum_{n\ge1}\frac1{16n^2}=\frac{\pi^2}{96}.
\tag{D6}\label{eq:D6}
\end{equation}
Each $\omega_s$, $s>0$, has finite support. The symbol $\omega_0$ denotes only
the $\ell^2$ comparison vector $(1/(4n))_{n\ge1}$ and is never an actual
graph parameter. Define
\begin{equation}
\begin{aligned}
J&=(-\infty,-1)\cup(1,\infty),\\
F(t)&=(\sgn t,\omega_{|t|-1})\in\ell^1(\Nzero)\quad(t\in J),\\
D&=\{a\in\ell^1(I):r(a)\in J,\ ma=F(r(a))\}.
\end{aligned}
\tag{D7}\label{eq:D7}
\end{equation}

The operators and kernel are
\begin{equation}
\begin{aligned}
\gra A&=\{(La,a):a\in D\},\\
K&=\{k\in\ell^1(I):mk=0,\ r(k)=0\},\\
Px&=\pair{x}{g}g\quad(x\in Y).
\end{aligned}
\tag{D8}\label{eq:D8}
\end{equation}
For each $t\in J$, the following is one element of $D$, not the entire fibre:
\begin{equation}
a^t=-t e_{0,1}+(t+\sgn t)e_{0,3}
       +\sum_{n\ge1}\omega_{|t|-1}(n)e_{n,1}.
\tag{D9}\label{eq:D9}
\end{equation}

\ifdefined\expandedversion
For standard $\ell^1$, let $Z=\ell^1(\N)$ and $Z^*=\ell^\infty(\N)$,
with their usual pairing. Fix once and for all an enumeration $(q_n)$ of
all finitely supported rational vectors in the closed unit ball of $Y$.
Define
\begin{equation}
\begin{aligned}
Qu&=\sum_{n\ge1}u_nq_n,& (Q^*a)_n&=\pair{q_n}{a},\\
\gra T&=\{(u,Q^*a):a\in D,\ Qu=La\},\\
f&=Q^*g,& Bu&=\pair{u}{f}f.
\end{aligned}
\tag{D10}\label{eq:D10}
\end{equation}
The proof establishes surjectivity with a norm-$2$ lift. No bounded linear
right inverse or complemented kernel is assumed. The polar for $T$ and
$T+B$ uses all $\ell^\infty$ labels.
\fi

\section{A counterexample on $c_0$}
\label{sec:c0}
\label{sec:statements}
\renewcommand{\thetheorem}{C}
\begin{theorem}[$c_0$ core]\label{thm:C}
Let $I=\Nzero\times\N$ and $Y=c_0(I)$ with its usual supremum norm and
full dual $Y^*=\ell^1(I)$. Let $A:Y\rightrightarrows Y^*$ and the bounded
positive rank-one map $P:Y\to Y^*$ be defined by Eq.~\eqref{eq:D8}.
Then:
\begin{enumerate}
\item $A$ and $P$ are maximally monotone in $Y\times Y^*$.
\item $\dom P=Y$ and $\dom A\ne\varnothing$, so
$\dom A\cap\operatorname{int}(\dom P)\ne\varnothing$.
\item Every $(x,a)\in\gra A$ satisfies $\|x\|_\infty>1/2$ and
$\pair{x}{a}\ge-2\sigma\|x\|_\infty$, where $\sigma=\pi^2/96<1/8$.
\item Every $(x,b)\in\gra(A+P)$ satisfies
$\pair{x}{b}>1-\sigma>7/8$.
\item $(0,0)$ is monotonically related to the entire graph of $A+P$ and
does not belong to that graph. Hence $A+P$ is not maximally monotone.
\end{enumerate}
\end{theorem}
\begin{proof}
\noindent\textit{Bounded maps and the bilinear identity.}
\label{sec:C1}
For $a\in\ell^1(I)$, $\|ma\|_1\le\|a\|_1$ and
$\|Ea\|_\infty\le2\|a\|_1$. If $a$ has finite support, then $Ea$ has
finite support: only finitely many blocks occur, and within a block no
output occurs after the last nonzero input coordinate. Approximation by
finite-support vectors and the bound prove $Ea\in c_0(I)$ globally.
Thus $L$ also maps $\ell^1(I)$ boundedly into $c_0(I)$.

For $a,c\in\ell^1(I)$, absolute convergence permits exchanging the double
sums. In each block the diagonal terms and the two triangular
off-diagonal sums give
\begin{equation}
\pair{Ea}{c}+\pair{Ec}{a}
 =2\sum_b\Big(\sum_j a_{b,j}\Big)\Big(\sum_j c_{b,j}\Big)
 =2\pair{ma}{mc}_{\ell^2}.
\tag{C1}\label{eq:C1}
\end{equation}
The absolute sum of all products is at most $\|a\|_1\|c\|_1$ before the
factor $2$, so this is an identity for all labels. Since $mg=0$, $h=Eg$,
and $r(g)=0$, Eq.~\eqref{eq:C1} implies $\pair{Ea}{g}=-r(a)$.
Consequently
\begin{equation}
\begin{aligned}
\pair{La}{a}&=r(a)^2-\|ma\|_2^2,&
\pair{La}{g}&=r(a),\\
\pair{La-Lc}{a-c}&=(r(a)-r(c))^2-\|ma-mc\|_2^2.
\end{aligned}
\tag{C2}\label{eq:C2}
\end{equation}

\noindent\textit{The Lipschitz curve and its complete fibres.}
\label{sec:C2}
For each $s>0$, $\omega_s(n)=0$ whenever $n\ge s^{-4}$. Also
$0\le\omega_s(n)\le1/(4n)$. Hence $W(s)\le\sigma<1/8$, and dominated
convergence in $\ell^2$ gives $\omega_s\to\omega_0$ and $W(s)\to\sigma$
as $s$ decreases to zero. The comparison vector $\omega_0$ is not in
$\ell^1$ by divergence of the harmonic series. No endpoint is added to $J$.

Since $v\mapsto\max\{v,0\}$ is $1$-Lipschitz on $\R$,
\begin{equation}
\|\omega_s-\omega_q\|_2^2\le c|s-q|^2,\qquad
c=\frac1{16}\sum_{n\ge1}n^{-3/2}<\frac3{16}<1.
\tag{C3}\label{eq:C3}
\end{equation}
The strict numerical bound follows from
$\sum_{n\ge1}n^{-3/2}<1+\int_1^\infty x^{-3/2}\,dx=3$.
For $t,u$ on the same branch of $J$, Eq.~\eqref{eq:C3} gives
$\|F(t)-F(u)\|_2\le|t-u|$. For $t=1+s$ and $u=-1-q$, the squared
distance is at most $4+c(s-q)^2\le(2+s+q)^2$, and the opposite
orientation is identical. Thus $F$ is $1$-Lipschitz on all of $J$.

The finitely supported vector $a^t$ in Eq.~\eqref{eq:D9} satisfies $r(a^t)=t$ and
$ma^t=F(t)$, proving nonemptiness for every actual $t$. The complete
fibre $\{a:r(a)=t,\ ma=F(t)\}$ is precisely $a^t+K$.
Applying Eq.~\eqref{eq:C2} and the Lipschitz inequality to any two labels
in $D$ proves monotonicity of the entire $A$.

\noindent\textit{Reduction of an arbitrary polar point.}
\label{sec:C3}
Take an arbitrary $(x,p)$ in the monotone polar, in the full pair
$Y\times Y^*$. Expansion using Eq.~\eqref{eq:C1} gives, for every
$a\in D$ and $t=r(a)$,
\begin{equation}
\begin{aligned}
\pair{x-La}{p-a}
 &=\pair{x}{p}-\pair{x+Ep}{a}\\
 &\quad+2\pair{ma}{mp}-t r(p)+t^2-\|ma\|_2^2.
\end{aligned}
\tag{C4}\label{eq:C4}
\end{equation}
Replace $a$ by $a+\theta k$ for $k\in K$ and all $\theta\in\R$.
All terms except $-\theta\pair{x+Ep}{k}$ are constant.
Nonnegativity for all $\theta$ forces $\pair{x+Ep}{k}=0$ for every $k\in K$.

We now prove $K^\perp\cap c_0(I)=\R h$ with no dual-space enlargement.
Let $z\in c_0(I)$ globally annihilate $K$. In a block $b\ge1$, every
difference $e_{b,j}-e_{b,k}$ is in $K$. Thus $z$ is constant in that
entire block, and global $c_0$ forces that constant to vanish. In block
zero, differences between indices $j,k\ge3$ belong to $K$, so the common
tail value is also zero. Finally, $g=e_{0,1}-e_{0,2}$ belongs to $K$,
forcing $z_{0,1}=z_{0,2}$. It follows that $z$ is a scalar multiple of
$h$. Conversely, $\pair{h}{k}=r(k)=0$ proves that every multiple of $h$
annihilates $K$. Therefore $x=-Ep+vh$ for some $v\in\R$.

Put $r_p=r(p)$, $m_p=mp$, $\tau=(v+r_p)/2$, and $\nu=(v-r_p)/2$.
Substitution into Eq.~\eqref{eq:C4}, including
$\pair{Ep}{p}=\|mp\|_2^2$, yields the exact bracket
\begin{equation}
\pair{x-La}{p-a}
 =(t-\tau)^2-\nu^2-\|F(t)-m_p\|_2^2.
\tag{C5}\label{eq:C5}
\end{equation}
Every $t\in J$ has an actual $a^t$, and Eq.~\eqref{eq:C5} is constant
over its complete fibre. Thus the full polar condition is exactly
\begin{equation}
\|F(t)-m_p\|_2^2+\nu^2\le(t-\tau)^2
\qquad\forall t\in J.
\tag{C6}\label{eq:C6}
\end{equation}

\noindent\textit{Exclusion of the missing parameter interval.}
\label{sec:C4}
If $|\tau|>1$, use the actual $t=\tau$ in Eq.~\eqref{eq:C6}. It forces
$\nu=0$ and $m_p=F(\tau)$. Then $v=r_p=\tau$, so $p\in D$ and $x=Lp$,
yielding graph membership.

If $|\tau|\le1$, write $m_p=(b,z_1,z_2,\ldots)$. For each fixed positive
integer $N$, discard the nonnegative tail of Eq.~\eqref{eq:C6}, then
take actual $t\to1$ from above and $t\to-1$ from below.
Finite-coordinate convergence gives
\begin{equation}
\begin{aligned}
(b-1)^2+\nu^2+\sum_{n=1}^N(z_n-1/(4n))^2&\le(1-\tau)^2,\\
(b+1)^2+\nu^2+\sum_{n=1}^N(z_n-1/(4n))^2&\le(1+\tau)^2.
\end{aligned}
\tag{C7}\label{eq:C7}
\end{equation}
Weight these by $(1+\tau)/2$ and $(1-\tau)/2$, respectively. Both
weights are nonnegative, including $\tau=1$ and $\tau=-1$. Expanding
the two squares gives
\begin{equation}
(b-\tau)^2+\nu^2+\sum_{n=1}^N(z_n-1/(4n))^2\le0.
\tag{C8}\label{eq:C8}
\end{equation}
For every $N$, all these nonnegative terms vanish. Thus $z_n=1/(4n)$ for
all $n$, contradicting $m_p\in\ell^1(\Nzero)$. No unattainable endpoint
has been used as a graph row, and no infinite interchange is required in
Eq.~\eqref{eq:C7}. This case is impossible. All polar points belong to
$A$, and monotonicity gives the reverse inclusion. Hence $A$ is maximally
monotone in $Y\times Y^*$.

\noindent\textit{The missing origin and the radial bound.}
\label{sec:C5}
For every $a\in D$, $t=r(a)$ satisfies $|t|>1$. By Eq.~\eqref{eq:C2}
and $\|g\|_1=2$,
\begin{equation}
\|La\|_\infty\ge|t|/2>1/2,\qquad
\pair{La}{a}=t^2-1-W(|t|-1)>-\sigma.
\tag{C9}\label{eq:C9}
\end{equation}
It follows that $A(0)=\varnothing$ and
$\pair{La}{a}\ge-2\sigma\|La\|_\infty$ for every actual graph label.
In particular $\mathcal V(\gra A)\le2\sigma<1/4$.
Indeed, for each $a\in D$, Eq.~\eqref{eq:C9} gives
$\max\{0,-\pair{La}{a}\}/\|La\|_\infty
\le\sigma/\|La\|_\infty<2\sigma$; taking the supremum proves the bound.

\noindent\textit{The rank-one factor and nonmaximality of the sum.}
\label{sec:C6}
The map $P$ is everywhere defined, linear, bounded with
$\|P\|\le\|g\|_1^2=4$, nonzero because $Pe_{0,1}=g$, and has range $\R g$.
Its monotonicity follows from $\pair{x-y}{Px-Py}=\pair{x-y}{g}^2$.

For completeness, suppose $(z,z^*)$ is related to its entire graph.
Testing at $z+\lambda v$ for any $v\in Y$ and all $\lambda\in\R$ gives
\begin{equation}
-\lambda\pair{v}{z^*-Pz}+\lambda^2\pair{v}{g}^2\ge0.
\tag{C10}\label{eq:C10}
\end{equation}
If the linear coefficient is nonzero, a sufficiently small $\lambda$ of
the opposite sign makes the expression negative. Hence $z^*=Pz$, proving maximality.

The actual label $a^2=-2e_{0,1}+3e_{0,3}$ gives
$La^2=-6e_{0,1}-8e_{0,2}-3e_{0,3}$. Thus
$\dom A\cap\operatorname{int}(\dom P)\ne\varnothing$, with
$\operatorname{int}(\dom P)=Y$. Every sum point has the form
$(La,a+tg)$, $a\in D$ and $t=r(a)$, and
\begin{equation}
\pair{-La}{-a-tg}=2t^2-1-W(|t|-1)>1-\sigma>7/8.
\tag{C11}\label{eq:C11}
\end{equation}
This proves that $(0,0)$ is related to every sum point. The sum is
monotone as the sum of two monotone operators; it has no zero primal by
Eq.~\eqref{eq:C9}. Adding $(0,0)$ is therefore a proper monotone extension.
This completes the proof of Theorem~\ref{thm:C}.
\end{proof}

\ifdefined\expandedversion
\section{The counterexample on standard $\ell^1$}
\label{sec:l1}
\renewcommand{\thetheorem}{L}
\begin{theorem}[Standard $\ell^1$ realization]\label{thm:L}
On $Z=\ell^1(\N)$ with its usual norm and full dual $Z^*=\ell^\infty(\N)$,
let $T:Z\rightrightarrows Z^*$ and $B:Z\to Z^*$ be defined by Eq.~\eqref{eq:D10}.
Then:
\begin{enumerate}
\item $T$ is maximally monotone in the full pair $Z\times Z^*$.
\item $B$ is nonzero, bounded, positive, rank-one, and maximally monotone,
with $\dom B=Z$. In particular,
$\dom T\cap\operatorname{int}(\dom B)\ne\varnothing$.
\item $T(0)=\varnothing$, and every $(u,b)\in\gra T$ satisfies
$\|u\|_1>1/2$ and $\pair{u}{b}\ge-2\sigma\|u\|_1$.
Thus $\mathcal V(\gra T)\le2\sigma<1/4$.
\item Every $(u,c)\in\gra(T+B)$ satisfies
$\pair{u}{c}>1-\sigma>7/8$.
\item $(0,0)$ belongs to the full monotone polar of $\gra(T+B)$ and does
not belong to $\gra(T+B)$. Hence $T+B$ is not maximally monotone.
\end{enumerate}
\end{theorem}
\begin{proof}
\noindent\textit{A surjection onto $c_0$.}
\label{sec:L1}
Finitely supported rational vectors in the closed unit ball of $Y$ are
countable and dense in that ball: truncate a $c_0$ vector and approximate
each retained coordinate by a rational inside $[-1,1]$. Fix an enumeration
of all these vectors as $(q_n)$. For $u\in Z=\ell^1(\N)$, the sum defining
$Qu$ converges absolutely in the Banach space $Y$, and
$\|Qu\|_\infty\le\|u\|_1$. Coordinate unit vectors occur among the $q_n$,
so $\|Q\|=1$.

For an arbitrary $x\in Y$, set $r_0=x$. If $r_k\ne0$, choose an actual
$q_{n_k}$ with $\|q_{n_k}-r_k/\|r_k\|_\infty\|_\infty<1/2$ and set
\begin{equation}
c_k=\|r_k\|_\infty,\qquad r_{k+1}=r_k-c_kq_{n_k}.
\tag{L1}\label{eq:L1}
\end{equation}
If a residual vanishes, stop. Otherwise
$\|r_k\|_\infty\le2^{-k}\|x\|_\infty$ and $\sum_k c_k\le2\|x\|_\infty$.
Define $u_n=\sum_{k:n_k=n}c_k$. Repeated indices are combined.
Positivity of the $c_k$ and summability give $u\in\ell^1$ with
$\|u\|_1=\sum_k c_k\le2\|x\|_\infty$.
Telescoping Eq.~\eqref{eq:L1}, then absolute convergence under regrouping,
proves
\begin{equation}
Qu=x,\qquad \|u\|_1\le2\|x\|_\infty.
\tag{L2}\label{eq:L2}
\end{equation}
For $x=0$, use $u=0$. This proves surjectivity with the asserted lifting
bound. The construction does not assert that $x\mapsto u$ is linear or
continuous, nor that $\ker Q$ is complemented.

\noindent\textit{The adjoint and arbitrary full-dual polar labels.}
\label{sec:L2}
Absolute convergence gives $\pair{u}{Q^*a}=\pair{Qu}{a}$ and
$\|Q^*a\|_\infty\le\|a\|_1$. For each finite subset of $I$, its sign
vector for $a$ is a finitely supported rational vector of norm at most
one and occurs in the enumeration. Taking these finite sign tests proves
\begin{equation}
\|Q^*a\|_\infty=\|a\|_1\qquad(a\in\ell^1(I)).
\tag{L3}\label{eq:L3}
\end{equation}
In particular, $Q^*$ is injective. The graph of $T$ is nonempty by
Eq.~\eqref{eq:L2} and nonemptiness of $D$. It is monotone because the
bracket of any two rows $(u,Q^*a)$ and $(v,Q^*c)$ equals
$\pair{La-Lc}{a-c}$, which is nonnegative by the $c_0$ core.

Let $(v,v^*)$ be a polar point with $v^*$ an arbitrary element of
$\ell^\infty(\N)$. Each actual fibre of $T$ contains $u+\lambda k$ for
every $k\in\ker Q$ and every $\lambda\in\R$. Testing those rows forces
$\pair{k}{v^*}=0$ for every such $k$. Define a functional on $Y$ by
\begin{equation}
\lambda(x)=\pair{u}{v^*}\quad\text{when }Qu=x.
\tag{L4}\label{eq:L4}
\end{equation}
Two preimages differ by $\ker Q$, proving well-definedness. Linear
combinations of preimages prove linearity. For each $x$, choose the
preimage in Eq.~\eqref{eq:L2}; then
$|\lambda(x)|\le2\|v^*\|_\infty\|x\|_\infty$. Thus $\lambda$ belongs to
the continuous dual of $c_0(I)$.

Explicitly put $p_i=\lambda(e_i)$. Finite sign tests give
$\sum_{i\in H}|p_i|\le\|\lambda\|$ for every finite $H$, whence
$p\in\ell^1(I)$. Equality $\lambda(x)=\pair{x}{p}$ holds first for
finite-support $x$ and then for every $c_0$ vector by density and continuity.
Therefore $v^*=Q^*p$: the two functionals agree on every $u$ because
Eq.~\eqref{eq:L4} applies to $Qu$.

Every original $a\in D$ admits a lift $u$ of $La$ by Eq.~\eqref{eq:L2}.
The full $T$-polar inequality now reads
\begin{equation}
\pair{Qv-La}{p-a}\ge0\qquad\forall a\in D.
\tag{L5}\label{eq:L5}
\end{equation}
The original maximality proof gives $p\in D$ and $Qv=Lp$. Hence
$(v,v^*)=(v,Q^*p)$ is an actual $T$ row. This proves maximality in the
full standard pair $\ell^1\times\ell^\infty$. Membership of $v^*$ in
$\operatorname{range}(Q^*)$ was derived from the actual kernel tests.

\noindent\textit{The radial bound on every graph value.}
\label{sec:L3}
For every $(u,Q^*a)\in\gra T$, with $t=r(a)$,
\begin{equation}
\begin{aligned}
\pair{u}{Q^*a}&=\pair{La}{a}
 =t^2-1-W(|t|-1)>-\sigma,\\
\|u\|_1&\ge\|Qu\|_\infty=\|La\|_\infty>1/2.
\end{aligned}
\tag{L6}\label{eq:L6}
\end{equation}
Thus $T(0)=\varnothing$ and $\pair{u}{Q^*a}\ge-2\sigma\|u\|_1$ on
the entire graph, including every $\ker Q$ translate. Consequently,
$\mathcal V(\gra T)\le2\sigma<1/4$. No exact equality for this
norm-dependent supremum is asserted.

\noindent\textit{The rank-one factor and the intersection condition.}
\label{sec:L4}
Let $f=Q^*g$. By Eq.~\eqref{eq:L3}, $\|f\|_\infty=\|g\|_1=2$, so $f$
is nonzero. The map $B:u\mapsto\pair{u}{f}f$ is nonzero, linear, bounded
with norm at most $4$, positive, and rank-one. The same whole-line proof
as Eq.~\eqref{eq:C10}, now testing all $v\in\ell^1$ and labels in
$\ell^\infty$, proves $B$ maximally monotone.
Its domain is all of $Z$, so $\operatorname{int}(\dom B)=Z$.

The vector $q_\circ=-(3/4)e_{0,1}-e_{0,2}-(3/8)e_{0,3}$ occurs as
$q_{n_0}$. For $u_\circ=8e_{n_0}$, $Qu_\circ=La^2$.
Thus $(u_\circ,Q^*a^2)$ is an actual $T$ row, and
$\dom T\cap\operatorname{int}(\dom B)\ne\varnothing$ in the stated
factor order.

\noindent\textit{A strict polar witness outside the graph.}
\label{sec:L5}
For every row $(u,Q^*a)$ of $T$, let $t=r(a)$. Then
$\pair{u}{f}=\pair{Qu}{g}=\pair{La}{g}=t$.
Consequently each sum row has work
\begin{equation}
\pair{u}{Q^*a+Bu}
 =2t^2-1-W(|t|-1)>1-\sigma>7/8.
\tag{L7}\label{eq:L7}
\end{equation}
The sum is monotone by monotonicity of its factors.
Equation~\eqref{eq:L7} shows that $(0,0)$ is related to all its rows.
Since $T(0)=\varnothing$, $(0,0)$ is not a sum row. Adjoining this point
gives a proper monotone extension, completing the proof of
Theorem~\ref{thm:L} with the standard Banach space, full dual, and ordered
intersection condition. No theorem about sums was used to prove either
factor maximal.
\end{proof}

\noindent\textit{A determined lift, not a choice of maximal extension.}
For a fixed enumeration $(q_n)$ in Eq.~\eqref{eq:D10}, the lift used in
Eq.~\eqref{eq:L2} can be specified without further choices. Given $x\in Y$,
put $x_0=x$. Whenever $x_k\ne0$, define
\[
n_k=\min\left\{n:\left\|q_n-\frac{x_k}{\|x_k\|_\infty}\right\|_\infty
 <\frac12\right\},\qquad
c_k=\|x_k\|_\infty,\qquad x_{k+1}=x_k-c_kq_{n_k},
\]
and stop if a residual is zero. Density makes the minimum well-defined,
and $\|x_k\|_\infty\le2^{-k}\|x\|_\infty$. Hence
\[
\mathcal R(x)=\sum_k c_ke_{n_k}\in\ell^1,\qquad
\|\mathcal R(x)\|_1\le2\|x\|_\infty,\qquad
Q\mathcal R(x)=x,
\]
where the last identity follows by telescoping and continuity of $Q$.
In particular,
\[
\gra T=\{(\mathcal R(La)+v,Q^*a):a\in D,\ v\in\ker Q\}.
\]
Thus the whole graph of $T$ is specified before maximality is proved;
no arbitrary maximal extension is selected. It is a quotient pullback of
the coordinate construction on $c_0$, not an independent coordinate
construction on $\ell^1$. Here ``specified'' does not assert constructivity
in the sense of constructive analysis or computability theory. In particular,
this assertion does not
claim a bounded linear or continuous right inverse, or an algorithm
deciding membership in the infinite-dimensional graph.
\fi

\section{An exact endpoint criterion}
\label{sec:endpoint-criterion}
The nonsummable tail has a precise role in the construction: the endpoint
comparisons force a dual label that the actual dual space cannot contain.
We now separate this requirement from the particular triangular coordinates.
The following theorem gives a criterion within a specified class of graphs;
it does not assume that arbitrary maximally monotone operators admit this
representation.

Let $X$ be a real Banach space and let $H_0$ be a real Hilbert space.
In this section put $H=\R\oplus H_0$. Let
$E:X^*\to X$ and $M:X^*\to H$ be bounded linear maps, and let
$g\in X^*\setminus\{0\}$ satisfy $Mg=0$ and
\begin{equation}
\pair{Ea}{b}+\pair{Eb}{a}=2(Ma,Mb)_H
\qquad(a,b\in X^*).
\label{eq:abs-bilinear}
\end{equation}
Define $h=Eg$, $r(a)=\pair{h}{a}$, $La=-Ea+r(a)h$, and
$K=\ker M\cap\ker r$. Assume the exact primal annihilator identity
\begin{equation}
\{x\in X:\pair{x}{k}=0\ \hbox{for all }k\in K\}=\R h.
\label{eq:abs-annihilator}
\end{equation}
Let $\omega:(0,\infty)\to H_0$ be $\ell$-Lipschitz, where
$0\le\ell\le1$, and suppose
\begin{equation}
\omega(s)\longrightarrow\eta\text{ in }H_0\quad(s\downarrow0),
\qquad 0<S:=\|\eta\|^2<1,\qquad \|\omega(s)\|^2\le S.
\label{eq:abs-tail}
\end{equation}
Put $J=(-\infty,-1)\cup(1,\infty)$ and
$F(t)=(\sgn t,\omega(|t|-1))$. For every $t\in J$, assume
there exists an actual $a^t\in X^*$ with $r(a^t)=t$ and
$Ma^t=F(t)$. We use every label with these values, defining
\begin{equation}
D=\{a\in X^*:r(a)\in J,\ Ma=F(r(a))\},\qquad
G=\{(La,a):a\in D\}.
\label{eq:abs-full-graph}
\end{equation}
In particular, each parameter fibre is $a^t+K$, without a boundedness
or continuity assumption on the choice $t\mapsto a^t$.

\renewcommand{\thetheorem}{E}
\begin{theorem}[Actual endpoint labels and maximality]
\label{thm:endpoint-criterion}
Under the preceding assumptions, define
\[
\widehat F(t)=
\begin{cases}
F(t),&|t|>1,\\
(t,\eta),&|t|\le1.
\end{cases}
\]
Then the entire monotone polar in $X\times X^*$ is
\begin{equation}
G^\mu=\{(Lp,p):p\in X^*,\ Mp=\widehat F(r(p))\}.
\label{eq:abs-entire-polar}
\end{equation}
It is the unique maximally monotone extension of $G$. Moreover, $G$
itself is maximally monotone if and only if
\begin{equation}
\nexists(p,\tau)\in X^*\times[-1,1]:
\qquad r(p)=\tau,\quad Mp=(\tau,\eta).
\label{eq:abs-endpoint-exclusion}
\end{equation}
When Eq.~\eqref{eq:abs-endpoint-exclusion} holds, the operator
$\mathcal A$ with graph $G$ omits zero and satisfies
$\mathcal V(G)\le S\|g\|$. With
$\mathcal P x=\pair{x}{g}g$, both $\mathcal A$ and $\mathcal P$
are maximally monotone, satisfy
$\dom\mathcal A\cap\operatorname{int}\dom\mathcal P\ne\varnothing$,
and have a nonmaximal sum: $(0,0)$ belongs to its monotone polar
and lies outside its graph.
\end{theorem}
\begin{proof}
Equation~\eqref{eq:abs-bilinear} at $a=b=g$ gives $r(g)=0$.
Using $Mg=0$ again gives
\begin{equation}
\pair{Ea}{g}=-r(a),\quad \pair{La}{g}=r(a),\quad
\pair{La-Lb}{a-b}=|r(a)-r(b)|^2-\|Ma-Mb\|^2.
\label{eq:abs-pairing}
\end{equation}
The actual label at $t=2$ also proves $h\ne0$.

Extend $\omega$ to zero by $\omega(0)=\eta$; its Lipschitz bound
persists by the norm limit. Define
$\chi(t)=\max\{-1,\min\{1,t\}\}$ and
$\zeta(t)=\max\{|t|-1,0\}$. Then
$\widehat F(t)=(\chi(t),\omega(\zeta(t)))$.
If $t,u$ lie on the same exterior branch, in the middle interval,
or one in the middle and one outside, the nonnegative changes in
$\chi$ and $\zeta$ add to $|t-u|$. On opposite exterior branches
they sum to $2+||t|-|u||\le|t-u|$. Thus in all cases
\[
|\chi(t)-\chi(u)|+|\zeta(t)-\zeta(u)|\le|t-u|,
\]
and
\[
\|\widehat F(t)-\widehat F(u)\|^2
\le|\chi(t)-\chi(u)|^2+\ell^2|\zeta(t)-\zeta(u)|^2
\le|t-u|^2.
\]
Equation~\eqref{eq:abs-pairing} proves monotonicity of both $G$
and the set on the right of Eq.~\eqref{eq:abs-entire-polar}.
It also proves that this latter set is contained in $G^\mu$.

For the reverse inclusion, take an arbitrary $(x,p)\in G^\mu$,
where $p$ belongs to the full continuous dual. For $a\in D$ and
$t=r(a)$, expansion using Eq.~\eqref{eq:abs-bilinear} gives
\begin{align*}
\pair{x-La}{p-a}
={}&\pair{x}{p}-\pair{x+Ep}{a}
       +2(Ma,Mp)-t r(p)+t^2-\|Ma\|^2.
\end{align*}
Replacing $a$ by $a+\theta k$, for all $\theta\in\R$ and
$k\in K$, forces $\pair{x+Ep}{k}=0$. Hence
$x=-Ep+v h$ by Eq.~\eqref{eq:abs-annihilator}. With
\[
\tau=\frac{v+r(p)}2,\qquad \nu=\frac{v-r(p)}2,
\]
the same expansion reduces exactly to
\begin{equation}
\pair{x-La}{p-a}=(t-\tau)^2-\nu^2-\|F(t)-Mp\|^2.
\label{eq:abs-polar-square}
\end{equation}
Every $t\in J$ is an actual test. If $|\tau|>1$, testing
$t=\tau$ forces $\nu=0$ and $Mp=F(\tau)$; thus
$v=r(p)=\tau$ and $x=Lp$.

If $|\tau|\le1$, write $Mp=(b,z)\in\R\oplus H_0$.
Limits along the two actual branches yield
\[
(b-1)^2+\|z-\eta\|^2+\nu^2\le(1-\tau)^2,
\qquad
(b+1)^2+\|z-\eta\|^2+\nu^2\le(1+\tau)^2.
\]
Multiply these by $(1+\tau)/2$ and $(1-\tau)/2$, respectively,
and add. Expansion gives
\[
(b-\tau)^2+\|z-\eta\|^2+\nu^2\le0.
\]
The endpoint weights remain valid at $\tau=\pm1$.
Consequently $Mp=(\tau,\eta)$, $r(p)=v=\tau$, and $x=Lp$.
This proves Eq.~\eqref{eq:abs-entire-polar}. The norm limits were
taken in $H$, not in $X^*$; the label $p$ was actual from the start.

We have proved that $G^\mu$ is monotone. Any point related to it
is related to $G$ and therefore belongs to $G^\mu$, so $G^\mu$
is maximally monotone. Every monotone extension of $G$ lies in
$G^\mu$; hence a maximal one must equal it. The missing part
$G^\mu\setminus G$ is exactly the set of labels in
Eq.~\eqref{eq:abs-endpoint-exclusion}. This proves the criterion.

Finally, each $a\in D$ gives
\[
|\pair{La}{g}|=|r(a)|>1,\quad
\|La\|>1/\|g\|,\quad
\pair{La}{a}=r(a)^2-1-\|\omega(|r(a)|-1)\|^2>-S.
\]
These imply the missing zero and the stated whole-graph radial bound.
The map $\mathcal P$ is bounded, linear, positive, nonzero, and full-domain.
For an arbitrary point $(z,z^*)$ in its polar, tests at $z+tv$ give
\[
-t\pair{v}{z^*-\mathcal Pz}+t^2\pair{v}{g}^2\ge0
\qquad(t\in\R,\ v\in X).
\]
Both signs of arbitrarily small $t$ force $z^*=\mathcal Pz$.
Thus $\mathcal P$ is maximally monotone, and the actual label at
$t=2$ gives the ordered intersection condition. For every sum row,
\[
\pair{La}{a+\mathcal PLa}
=2r(a)^2-1-\|\omega(|r(a)|-1)\|^2>1-S>0.
\]
The sum is monotone and omits the zero primal. Adjoining $(0,0)$
is a proper monotone extension.
\end{proof}

For the concrete example, use the maps $E,m,g$ of
Eqs.~\eqref{eq:D4}--\eqref{eq:D5}, $H_0=\ell^2(\N)$, and
$\omega(s)=\omega_s$. Equation~\eqref{eq:C1} supplies the
bilinear identity. The complete-kernel calculation in the proof of
Theorem~\ref{thm:C} supplies the primal annihilator identity,
and Eq.~\eqref{eq:D9} supplies every actual parameter fibre.
The tail bound and limit give $S=\sigma$. Since every $mp$ is
summable and $\eta=(1/(4n))_{n\ge1}$ is not, the endpoint
exclusion holds. Thus the theorem explains which properties of the
coordinate construction are needed and why the Lipschitz extension
$\widehat F$ in the Hilbert comparison space does not add a graph point.

\noindent\textit{Reverse design from the polar condition.}
The criterion gives a way to reconstruct the design requirements from
the full polar: retain all kernel fibres so that every polar query has
the scalar form in Eq.~\eqref{eq:abs-polar-square}, realize the two
exterior branches, and choose a common limiting tail with no actual
dual label. The bound in Eq.~\eqref{eq:abs-tail} then supplies the
finite radial estimate and the rank-one partner. For the present model,
the gap between $\ell^1$ and $\ell^2$ provides exactly that absent
label. This is a mathematical reconstruction of the mechanism, rather
than an additional assumption or a replacement for the preceding
coordinate proof.


\ifdefined\expandedversion
\section{A nonconstructive family with a different first factor}
\label{sec:nonconstructive}
The graphs in Theorems~\ref{thm:C} and~\ref{thm:L} are determined before
their maximality is proved. We next give a different construction in which
the first factor is selected by Zorn's lemma. Its premises are verified
by a countable set of specified rows. An additional row ensures that the
selected factor differs from the corresponding determined factor, for
every allowed selection. The endpoint obstruction remains the same;
this is not an independent mechanism for the sum question.

\begin{lemma}[An extension preserving a lower pairing bound]
\label{lem:capped-extension}
Let $X$ be a real Banach space and let $S\subset X\times X^*$ be a
nonempty monotone relation. Suppose $c,\rho>0$ satisfy
\[
\pair{x}{a}\ge-c\quad((x,a)\in S),\qquad
S^\mu\cap\bigl(\overline B_X(0,\rho)\times X^*\bigr)=\varnothing.
\]
There is a maximally monotone graph $G\supset S$ such that
\[
\pair{x}{a}\ge-c,\qquad \|x\|>\rho\quad((x,a)\in G).
\]
In particular, $0\notin\dom G$ and $\mathcal V(G)\le c/\rho$.
\end{lemma}
\begin{proof}
Order by inclusion the monotone relations containing $S$ whose every row
satisfies the lower pairing bound $-c$. The relation $S$ belongs to this
family. The union of a chain is monotone, since any two of its rows occur
together in one chain member; it also contains $S$ and preserves the
bound. Zorn's lemma supplies an inclusion-maximal member $G$.

We prove maximality in the full pair $X\times X^*$, not merely in this
ordered family. For pairs $u=(x,a)$ and $v=(y,b)$, write
$[u,v]=\pair{x-y}{a-b}$, and let $o=(0,0)$. If
$w\in G^\mu\setminus G$ and $[o,w]\ge-c$, adjoining $w$ contradicts
the choice of $G$. Otherwise put $d=-[o,w]>c$ and $\alpha=c/d\in(0,1)$.
For every $v\in G$, expansion of the pairing gives
\begin{align*}
[\alpha w,v]
&=(1-\alpha)[o,v]+\alpha[w,v]-\alpha(1-\alpha)[o,w]\\
&\ge-(1-\alpha)c+\alpha(1-\alpha)d=0.
\end{align*}
Thus $\alpha w\in G^\mu$, and
$[o,\alpha w]=-c^2/d>-c$. But
$[w,\alpha w]=-(1-\alpha)^2d<0$, so $\alpha w\notin G$ because
$w$ is related to every point of $G$. Adjoining $\alpha w$ is another
contradiction. Therefore $G^\mu=G$.

Since $S\subset G$, polar reversal gives $G=G^\mu\subset S^\mu$.
The assumed exclusion implies $\|x\|>\rho$ on every row. The lower
pairing bound then gives
$\max\{0,-\pair{x}{a}\}/\|x\|\le c/\rho$. Taking the supremum
proves the last assertion.
\end{proof}

\begin{lemma}[A full-domain norm partner]\label{lem:norm-partner}
Let $M:X\rightrightarrows X^*$ be maximally monotone, with
$0\notin\dom M$, on a real Banach space. For $\lambda>0$, let
$N_\lambda=\partial(\lambda\|\cdot\|)$. This operator is maximally
monotone and has full domain. If $\mathcal V(\gra M)\le\lambda$, then
$(M,N_\lambda)$ satisfies the interior-domain condition and
$M+N_\lambda$ is not maximally monotone, with polar witness $(0,0)$.
\end{lemma}
\begin{proof}
The subgradient inequality, tested at zero and at scalar multiples of
$x$, gives $\pair{x}{j}=\lambda\|x\|$. Substitution back into that
inequality gives $\pair{y}{j}\le\lambda\|y\|$ for every $y\in X$,
which is equivalent to $\|j\|\le\lambda$. Thus
\[
N_\lambda x=\{j\in X^*: \|j\|\le\lambda,
                  \pair{x}{j}=\lambda\|x\|\}.
\]
Conversely, these two conditions imply the subgradient inequality by the
dual norm bound. Hahn--Banach provides a norming functional for every
$x\ne0$, and $N_\lambda0=\lambda\overline B_{X^*}$. Hence the domain is
all of $X$. For $j\in N_\lambda x$ and $k\in N_\lambda y$,
\[
\pair{x-y}{j-k}
=\lambda\|x\|+\lambda\|y\|-\pair{x}{k}-\pair{y}{j}\ge0,
\]
so the graph is monotone.

Suppose $(z,p)$ is related to this entire graph. For any $v\in X$,
choose $j\in N_\lambda(z+v)$. The test at $z+v$ gives
$\pair{v}{p}\le\pair{v}{j}\le\lambda\|v\|$. Replacing $v$ by $-v$
shows $\|p\|\le\lambda$. If $z\ne0$, testing at $tz$ for $t>0$ yields
\[
(1-t)\bigl(\pair{z}{p}-\lambda\|z\|\bigr)\ge0.
\]
Values on both sides of $t=1$ force equality. At $z=0$ that equality is
automatic. Thus $p\in N_\lambda z$, proving maximality directly.
This is the norm case of the subdifferential maximality theorem
\cite{rockafellar1970subdiff}.

For every $(x,a)\in\gra M$ and every $j\in N_\lambda x$,
\[
\pair{x}{a+j}=\pair{x}{a}+\lambda\|x\|\ge0.
\]
Thus $(0,0)$ is related to every sum row. It is absent from the graph
because $M0=\varnothing$. The sum is monotone, and adjoining that point
gives a proper monotone extension. Finally, a maximally monotone graph is
nonempty, since the empty graph admits a one-point extension. Together
with $\dom N_\lambda=X$, this proves the ordered intersection condition.
\end{proof}

We now use the concrete maps of Eqs.~\eqref{eq:D4}--\eqref{eq:D9}.
Let $K_{\mathbb Q}$ be the set of finitely supported rational labels
$k\in\ell^1(I)$ satisfying $mk=0$ and $r(k)=0$. Put
\[
t_n^+=1+1/n,\quad t_n^-=-1-1/n\quad(n\ge1),\qquad
S_0=\{(La,a):a=a^{t_n^\pm}+k,\ n\ge1,\ k\in K_{\mathbb Q}\}.
\]
Define one more label and the augmented relation by
\begin{equation}
p_\circ=-4e_{0,1}+4e_{0,3},\qquad x_\circ=Lp_\circ,
\qquad S=S_0\cup\{(x_\circ,p_\circ)\}.
\label{eq:nonconstructive-seed}
\end{equation}

\renewcommand{\thetheorem}{N}
\begin{theorem}[A selected counterexample on both spaces]
\label{thm:nonconstructive}
The relation $S$ is a countable monotone subset of
$c_0(I)\times\ell^1(I)$. There exists a maximally monotone
$\mathcal A$ whose graph contains $S$ and whose every row satisfies
\[
\|x\|_\infty>1/2,\qquad \pair{x}{a}\ge-\sigma,
\qquad \mathcal V(\gra\mathcal A)\le2\sigma.
\]
Every such extension satisfies $\mathcal A\ne A$, and
$\mathcal A+\partial(2\sigma\|\cdot\|_\infty)$ is nonmaximal under
the original interior-domain condition.
With the specified $Q:\ell^1\to c_0(I)$, the operator
$\mathcal T=Q^*\mathcal A Q$ is maximally monotone in
$\ell^1\times\ell^\infty$, differs from $T$, and satisfies
\[
0\notin\dom\mathcal T,\qquad
\mathcal V(\gra\mathcal T)\le2\sigma.
\]
The pair $(\mathcal T,\partial(2\sigma\|\cdot\|_1))$ also satisfies
the original condition and has a nonmaximal sum. Both sums have
$(0,0)$ in their full monotone polar and outside their graph.
\end{theorem}
\begin{proof}
Finite rational labels form a countable set. Each $a^{t_n^\pm}$ has
finite support, so $S_0$ is countable. It is a subset of $\gra A$,
because adding $k\in K_{\mathbb Q}\subset K$ leaves $r$ and $m$
unchanged. Thus $S_0$ is nonempty and monotone. Its every row has
pairing greater than $-\sigma$ by Eq.~\eqref{eq:C9}.
The added label satisfies
\[
r(p_\circ)=4,\qquad mp_\circ=0,\qquad
x_\circ=-8e_{0,1}-12e_{0,2}-4e_{0,3},\qquad
\pair{x_\circ}{p_\circ}=16.
\]
For a row $(La,a)\in S_0$, put $t=r(a)=t_n^\pm$. Since $|t|\le2$,
Eq.~\eqref{eq:C2}, valid for all dual labels, gives
\[
\pair{x_\circ-La}{p_\circ-a}
=(4-t)^2-\|F(t)\|_2^2\ge3-\sigma>0.
\]
Consequently $S$ is monotone and obeys the lower pairing bound
$-\sigma$. Also $(x_\circ,p_\circ)\notin\gra A$, since
$mp_\circ=0\ne F(4)=(1,0,\ldots)$.

It remains to prove the full-polar exclusion required by
Lemma~\ref{lem:capped-extension}. Take any
$(x,p)\in S_0^\mu$, with $x\in c_0(I)$ and $p\in\ell^1(I)$.
Fixing any $t_n^\pm$ and testing the labels
$a^{t_n^\pm}+j k$ for all integers $j$ and all $k\in K_{\mathbb Q}$
in Eq.~\eqref{eq:C4} forces $\pair{x+Ep}{k}=0$.
All coordinate differences used in the annihilator calculation in the
proof of Theorem~\ref{thm:C} belong to $K_{\mathbb Q}$: they are
$e_{b,j}-e_{b,k}$ for $b\ge1$, those with $j,k\ge3$ in block zero,
and $g=e_{0,1}-e_{0,2}$. Their tests and the global $c_0$ condition
therefore give $x=-Ep+v h$ for some $v\in\R$.
In particular, $\pair{x}{g}=r(p)$ by Eq.~\eqref{eq:C1} and $r(g)=0$.

Suppose $\rho:=r(p)$ satisfies $|\rho|\le1$, and write $mp=(b,z)$.
For the actual tests $t=t_n^\pm$, Eq.~\eqref{eq:C4} becomes
\[
v\rho-(v+\rho)t+t^2-\|F(t)-mp\|_2^2\ge0.
\]
Use $\omega_{1/n}\to\omega_0$ in $\ell^2$, and put
$\eta=\omega_0$. Letting $n\to\infty$ on the two branches gives
\begin{align*}
v(\rho-1)-\rho+2b-b^2-\|z-\eta\|_2^2&\ge0,\\
v(\rho+1)+\rho-2b-b^2-\|z-\eta\|_2^2&\ge0.
\end{align*}
Multiply by $(1+\rho)/2$ and $(1-\rho)/2$, respectively, and add.
These are nonnegative weights even at $\rho=\pm1$; the terms in $v$
cancel, leaving
\[
-(b-\rho)^2-\|z-\eta\|_2^2\ge0.
\]
This forces $z=\eta$, which is impossible because $mp\in\ell^1$
whereas $\eta\notin\ell^1$. Thus every point of $S_0^\mu$ has
$|\pair{x}{g}|>1$ and $\|x\|_\infty>1/2$, since $\|g\|_1=2$.
Polar reversal gives the same exclusion for $S^\mu\subset S_0^\mu$.
All limits were comparisons with actual seed rows, not additions of
endpoint labels to the dual space.

Apply Lemma~\ref{lem:capped-extension} with $c=\sigma$ and radius
$1/2$. It produces $\mathcal A$ with the asserted maximality and
bounds. The forced row $(x_\circ,p_\circ)$ proves
$\mathcal A\ne A$ for every such choice. Lemma~\ref{lem:norm-partner}
proves the claimed counterexample on $c_0$. More explicitly, for every
$(x,a)\in\gra\mathcal A$ and every norm-subgradient value $j$,
\[
\pair{x}{a+j}\ge-\sigma+2\sigma\|x\|_\infty>0.
\]

For completeness, define the entire pullback graph by
\[
\gra\mathcal T=\{(u,Q^*a):(Qu,a)\in\gra\mathcal A\}.
\]
The surjectivity and norm-$2$ lifting bound of $Q$ were proved in
Eqs.~\eqref{eq:L1}--\eqref{eq:L2}, and the full-dual adjoint isometry
in Eq.~\eqref{eq:L3}; they do not depend on which monotone graph is
pulled back. This graph is nonempty and monotone by the adjoint identity.
For any polar query $(v,v^*)\in\ell^1\times\ell^\infty$, all translates
$u+t k$, $k\in\ker Q$, of one actual graph row force
$\pair{k}{v^*}=0$. The functional
$x\mapsto\pair{u}{v^*}$, $Qu=x$, is therefore well-defined and linear.
Choosing a norm-$2$ lift bounds it by $2\|v^*\|_\infty\|x\|_\infty$.
The full dual of $c_0(I)$ supplies $p\in\ell^1(I)$ representing this
functional, so $v^*=Q^*p$. Every row of $\mathcal A$ has a lift, and
the remaining polar inequalities are
\[
\pair{Qv-x}{p-a}\ge0\qquad((x,a)\in\gra\mathcal A).
\]
Maximality of $\mathcal A$ gives $(Qv,p)\in\gra\mathcal A$ and hence
$(v,v^*)\in\gra\mathcal T$. This proves full-dual maximality.

For every pullback row, $\|u\|_1\ge\|Qu\|_\infty>1/2$ and
$\pair{u}{Q^*a}=\pair{Qu}{a}\ge-\sigma$. Thus the asserted missing
origin and radial bound hold. A lift $u_\circ$ of $x_\circ$ gives a row
$(u_\circ,Q^*p_\circ)$ of $\mathcal T$. If it were a row of $T$,
injectivity of $Q^*$ would force its original label to be $p_\circ$,
contradicting $p_\circ\notin D$. Therefore $\mathcal T\ne T$.
Lemma~\ref{lem:norm-partner} now proves the second sum claim, including
full domain of its norm partner, all graph values, the ordered
intersection condition, and graph nonmembership of $(0,0)$.
\end{proof}

This theorem selects an unspecified member of a nonempty family. We do
not claim that two selections give different operators or that membership
in the selected graph is decidable. The forced row proves exactly the
distinction stated in the theorem. The proof also explains why a lower
pairing bound on seed rows alone is insufficient: the radial argument
is needed to promote a bound-preserving extension to global maximality.


\section{Every positive perturbation strength gives a counterexample}
\label{sec:parameters}
The preceding theorems use coefficient one. The same first operators
also give counterexamples for arbitrarily small positive coefficients.
The distinction between the origin and a general polar query is essential
for this strengthening.

\begin{proposition}\label{prop:all-strengths}
For the fixed operators of Theorems~\ref{thm:C} and \ref{thm:L}, every
$\lambda>0$ gives two maximally monotone factors satisfying the ordered
interior-domain condition, while neither $A+\lambda P$ nor $T+\lambda B$
is maximally monotone. More precisely,
\[
\{\lambda\in\R:A+\lambda P\text{ is maximally monotone}\}
=\{\lambda\in\R:T+\lambda B\text{ is maximally monotone}\}
=\{0\}.
\]
For $\lambda\ge0$, the origin belongs to the full monotone polar of
either sum exactly when $\lambda\ge\sigma$.
\end{proposition}
\begin{proof}
Fix $\lambda>0$. The whole-line argument in Eq.~\eqref{eq:C10}, with
the quadratic term multiplied by $\lambda$, proves maximality of
$\lambda P$ and $\lambda B$. They have full domain, so the same nonempty
domains of $A$ and $T$ verify the ordered intersection condition.

Set $s=\sqrt\lambda/2$ and define
\begin{equation}
p_s=\frac{a^{1+s}+a^{-1-s}}2
   =\sum_{n\ge1}\omega_s(n)e_{n,1},\qquad x_s=-Ep_s=Lp_s.
\label{eq:parameter-query}
\end{equation}
The sums are finite for each $s>0$. We have $r(p_s)=0$,
$mp_s=(0,\omega_s)$ and $\pair{x_s}{g}=0$. These are query coordinates,
not an assertion that $p_s\in D$.
For any $a\in D$, put $t=r(a)$ and $u=|t|-1>0$.
The identity in Eq.~\eqref{eq:C2} applies to arbitrary dual labels,
including $p_s$, and gives
\begin{equation}
\begin{aligned}
\pair{x_s-La}{p_s-a-\lambda t g}
 &=(1+\lambda)(1+u)^2-1-\|\omega_u-\omega_s\|_2^2\\
 &\ge\lambda-cs^2+2(1+\lambda+cs)u+(1+\lambda-c)u^2\\
 &>3\lambda/4>0,
\end{aligned}
\label{eq:parameter-bracket}
\end{equation}
where $c<1$ is the constant in Eq.~\eqref{eq:C3}.
Every primal in $\dom A$ has $|\pair{La}{g}|>1$, so $x_s\notin\dom A$.
Thus $(x_s,p_s)$ is related to every value of the sum graph and lies
outside it. Since the sum is monotone, adjoining this point gives a
proper monotone extension.

For the standard $\ell^1$ realization, choose $v_s$ with $Qv_s=x_s$
using Eq.~\eqref{eq:L2}. For every sum row $(u,Q^*a+\lambda t f)$,
the adjoint identity gives
\begin{equation}
\pair{v_s-u}{Q^*p_s-Q^*a-\lambda t f}
=\pair{x_s-La}{p_s-a-\lambda t g}>3\lambda/4.
\label{eq:parameter-l1}
\end{equation}
This covers all graph values and all $\ker Q$ translates.
Since $Qv_s=x_s\notin\dom A$, $v_s\notin\dom T$.
Consequently $(v_s,Q^*p_s)$ is a full-dual polar point outside the sum
graph. No bounded linear right inverse is used.

At $\lambda=0$, both first operators are maximal by the preceding
theorems. For $\lambda<0$, set $s=-\lambda/6>0$ and take the two actual
labels $a^{1+s}$ and $a^{-1-s}$. Their tails coincide. Equation~\eqref{eq:C2}
therefore gives the bracket of the corresponding perturbed graph rows as
\begin{equation}
4\big((1+\lambda)(1+s)^2-1\big)
=\frac{\lambda(\lambda-3)(\lambda-8)}9<0.
\label{eq:parameter-negative}
\end{equation}
Lifting their two primals through $Q$ gives the same negative bracket
for $T+\lambda B$. Neither sum is monotone at these parameters.
Here the negative multiples are used only to classify the parameter
sets; they are not claimed to be admissible monotone second factors.

Finally, for $\lambda\ge0$, every sum row in either realization has
self-pairing
$(1+\lambda)t^2-1-W(|t|-1)$, whose infimum over actual $|t|>1$ is
$\lambda-\sigma$. The lower bound follows from $W\le\sigma$ and
$t^2>1$; actual parameters tending to either deleted endpoint prove
the reverse bound. Thus the origin is polar exactly when
$\lambda\ge\sigma$. For $0<\lambda<\sigma$, the different query in
Eq.~\eqref{eq:parameter-query} remains valid.
\end{proof}
\fi

\section{Domain geometry and an almost-convexity statement}
\label{sec:domain-comparison}
The same graph also has a nonconvex domain closure. This consequence
helps identify the relationship between the construction and claims
about the domains of arbitrary maximally monotone operators.

\begin{proposition}\label{prop:domain-gap}
For the operator $A$ of Theorem~\ref{thm:C},
$0\in\operatorname{conv}(\dom A)$ and
$\operatorname{dist}(0,\dom A)\ge1/2$.
In particular, $\overline{\dom A}$ is not convex.
\end{proposition}
\begin{proof}
At $t=2$ and $t=-2$, the tails in Eq.~\eqref{eq:D9} vanish, so that
$a^{-2}=-a^2$. Since $L$ is linear, both $La^2$ and $-La^2$ belong
to $\dom A$. Their midpoint is zero. Equation~\eqref{eq:C9} gives
$\|x\|_\infty>1/2$ for every $x\in\dom A$, and hence the asserted
distance bound. A convex closure containing $La^2$ and $-La^2$ would
contain zero, contradicting this bound.
\end{proof}

Eberhard and Wenczel \cite[Theorem~57]{eberhardwenczel2024} state that
the norm closure of the domain of every maximally monotone operator on
a Banach space is convex. Proposition~\ref{prop:domain-gap}, together
with the proof of Theorem~\ref{thm:C}, gives an incompatible
conclusion for the displayed operator. The proof on
\cite[pp.~39--40]{eberhardwenczel2024} uses a small-parameter estimate
from Proposition~56, Eq.~(65), along a family in which the operator and
the evaluated point vary. The estimate for a fixed operator and point
has an object-dependent threshold; it does not supply a threshold
uniform over that family. The additional uniform estimate is needed
for that application. Our construction and maximality argument do not
use this published domain-closure statement or its negation.

There is also a version-specific issue with the separable-Asplund sum
statement in Bo\c{t} and Csetnek~\cite[Corollary~11]{botcsetnek2008}.
Its written hypotheses apply to the $c_0(I)$ construction: this space
is separable Asplund and $\dom P=Y$, so the convexified domain difference
is all of $Y$. The inspected proof invokes their Theorem~3, which asserts
bidual conjugate domination for representatives of every maximally
monotone operator when the dual space is separable. The operator in
\cite[author version, Lemma~2.1]{buenosvaiter2011} shows why that assertion cannot be used.
Write $G$ for the Gossez skew map on $\ell^1$, $e=(1,1,\ldots)$, and
$s(a)=\sum_n a_n$. That operator has graph
$\{(-Ga,a):a\in\ell^1,\ s(a)=0\}$. Its Fitzpatrick function is the
indicator of $x+Gb\in\R e$; since $x\in c_0$ and $Gb$ tends to $-s(b)$,
the finite-value points satisfy $x=-Gb-s(b)e$. For its conjugate on
$\ell^1\times\ell^\infty$, direct substitution therefore gives
\begin{equation}
\begin{aligned}
F_S^{c*}(p,\xi)
 &=\sup_{b\in\ell^1}\pair{Gp-s(p)e+\xi}{b},\\
F_S^{c*}(e_1,e-Ge_1)&=0<1=\pair{e-Ge_1}{e_1}.
\end{aligned}
\label{eq:bidual-interface}
\end{equation}
This calculation refutes the domination assertion used in that proof;
failure of the proof alone does not decide the standalone sum statement.
Theorem~\ref{thm:C} is proved independently of it.
\ifdefined\expandedversion
Although standard
$\ell^1$ is not Asplund, Theorem~\ref{thm:L} still uses the maximality
of $A$ established in Theorem~\ref{thm:C}; passing to $\ell^1$ does not
remove that dependency. The distinction concerns the cited 2008 preprint,
not an assertion about every later version or possible repair.
\else
The distinction concerns the cited 2008 preprint, not an assertion about
every later version or possible repair.
\fi

\ifdefined\expandedversion
\section{The structure shared by the examples}
\label{sec:structure}
The examples use one obstruction in two different ways. For the determined
graphs, it proves maximality by computing the entire monotone polar. For
the selected graphs, it first excludes an entire primal region from the
polar of a countable seed; the extension argument then supplies a maximal
graph with a lower pairing bound. The latter argument adds a choice of
extension, not a new kind of endpoint obstruction.

\subsection{Actual labels and the completed profile}
The coordinates $r(a)$ and $ma$ preserve the symmetric part of the graph
pairing:
\[
\pair{La-Lb}{a-b}=|r(a)-r(b)|^2-\|ma-mb\|_2^2.
\]
Consequently, a Lipschitz curve in the $(r,m)$ coordinates supplies
monotonicity. The complete kernel fibres supply different information:
they constrain every polar query to the same scalar comparison, rather
than merely proving inequalities between the chosen parameter rows.
Both ingredients are necessary in the proofs above. A parameter curve
alone is not an assertion of maximality for its image in $X\times X^*$.

In the architecture of Theorem~\ref{thm:endpoint-criterion}, the precise
criterion is Eq.~\eqref{eq:abs-endpoint-exclusion}. The comparison space
$\R\oplus H_0$ contains the middle segment $(t,\eta)$, $|t|\le1$,
but a point on that segment gives a graph point only if an actual
$p\in X^*$ realizes both $r(p)=t$ and $Mp=(t,\eta)$. The theorem proves
this statement in both directions. In the concrete construction,
$mp\in\ell^1$ rules out that realization. Thus passing from actual
block sums to their $\ell^2$ limits preserves the comparison geometry
but not the existence of dual labels. No assertion that $\ell^2$ is
the continuous dual of the original primal space is involved.

The following consequence makes the loss of dual control visible in
the original Banach spaces, without relying on a picture in the
comparison space.

\begin{proposition}[Convergent primals and unbounded dual labels]
\label{prop:boundary-sequences}
There are sequences of rows common to $A$ and every selected
$\mathcal A$ of Theorem~\ref{thm:nonconstructive} whose primal
coordinates converge in $c_0(I)$ to points outside both domains,
while the norms of their dual coordinates tend to infinity.
There are also such sequences common to $T$ and the corresponding
$\mathcal T$ on standard $\ell^1$, with convergence in its norm
and divergence in the full $\ell^\infty$ dual norm.
\end{proposition}
\begin{proof}
Use $a_n^\pm=a^{\pm(1+1/n)}$ and $x_n^\pm=La_n^\pm$. These rows
belong to $S_0$, hence to both first graphs on $c_0$.
Writing $t=\pm(1+1/n)$ and $\epsilon=\sgn t$, direct substitution in
Eqs.~\eqref{eq:D4}--\eqref{eq:D9} gives the block-zero coordinates
\[
(La^t)_{0,1}=-2t-2\epsilon,\qquad
(La^t)_{0,2}=-3t-2\epsilon,\qquad
(La^t)_{0,3}=-t-\epsilon,
\]
with every later coordinate in that block zero. In each block $b\ge1$,
only the first primal coordinate is nonzero, and it is
$-\omega_{1/n}(b)$. The convergence in $\ell^2$ of these tails implies
convergence in the supremum norm. The two limits $x_\pm$ therefore have
block-zero coordinates $\mp(4,5,2)$ and positive-block first coordinates
$-1/(4b)$. They belong to global $c_0(I)$, and
$\pair{x_\pm}{g}=\pm1$. Neither limit belongs to $\dom A$ or
$\dom\mathcal A$: the proofs of Theorems~\ref{thm:C} and
\ref{thm:nonconstructive} exclude every primal with
$|\pair{x}{g}|\le1$.

On the other hand,
\[
\|a_n^\pm\|_1=2(1+1/n)+1+\|\omega_{1/n}\|_1\longrightarrow\infty.
\]
To verify the divergence, fix any $N$. The sum of the first $N$ tail
coordinates tends to $\sum_{b=1}^N1/(4b)$. These partial harmonic sums
are unbounded as $N\to\infty$, so the full nonnegative tail sums tend
to infinity. This argument does not exchange two infinite sums.

For each sign choose a lift $u_\pm$ of $x_\pm$ through $Q$.
The norm-$2$ lifting bound gives $v_n^\pm$ with
$Qv_n^\pm=x_n^\pm-x_\pm$ and
$\|v_n^\pm\|_1\le2\|x_n^\pm-x_\pm\|_\infty$.
Then $u_n^\pm=u_\pm+v_n^\pm\to u_\pm$ in $\ell^1$, and
$(u_n^\pm,Q^*a_n^\pm)$ belongs to both pullback graphs. The limits
are outside both domains because their $Q$ images are $x_\pm$.
Equation~\eqref{eq:L3} gives
$\|Q^*a_n^\pm\|_\infty=\|a_n^\pm\|_1\to\infty$.
Only controlled individual lifts were used, not a continuous or linear
right inverse.
\end{proof}

\subsection{A lower pairing bound and different second factors}
The examples separate boundedness of the dual labels from control of
their pairing with their own primals. Proposition~\ref{prop:boundary-sequences}
shows that dual norms can diverge along convergent primals. Nevertheless,
the determined graph has the exact identity
\[
\pair{La}{a}=r(a)^2-1-W(|r(a)|-1)>-\sigma,
\]
and the selected graph retains the lower bound $-\sigma$ on every row.
Together with $\|x\|>1/2$, this proves the finite radial bound; it is
not a premise inferred from a conjectured failure of a sum theorem.

The second factors use that control in different ways. A rank-one factor
adds $\lambda r(a)^2$ on the determined graph. For $\lambda\ge\sigma$
the origin is a polar point of the sum. For $0<\lambda<\sigma$, the
finite-support query in Eq.~\eqref{eq:parameter-query} replaces it;
the full computation in Proposition~\ref{prop:all-strengths} is needed
at those strengths. A norm subdifferential instead adds exactly
$\lambda\|x\|$ for every one of its values, so that the lower bound
and the missing primal alone suffice. This is why it also works for
the unspecified first factor in Theorem~\ref{thm:nonconstructive}.
\ifdefined\expandedversion
The bounded-domain normal-cone example in Corollary~\ref{cor:normal}
uses the closed convex set $C_\rho$, which contains $0$. Testing its
normal inequality at $0$ gives $\pair{x}{n}\ge0$ for every
$x\in C_\rho$ and every $n\in N_{C_\rho}(x)$. The explicit interior
witness there verifies the ordered domain condition. That example is a change
of second factor, not a new independent construction of the first.
\fi

The extension proof identifies the additional requirement behind the
nonconstructive family. A lower bound on a seed is not automatically a
lower bound on its full polar, and inclusion-maximality among graphs
obeying that bound is not automatically global maximality. The scaled
polar point in Lemma~\ref{lem:capped-extension} closes precisely this
gap. At the same time, the seed excludes all dual values above the
forbidden primal slab, so every resulting maximal extension still omits
zero. The added row in Eq.~\eqref{eq:nonconstructive-seed} forces the
new factor to differ from the original one while preserving both
properties.

These distinctions delimit the representation result as well as its
use. Theorem~\ref{thm:endpoint-criterion} characterizes maximality for
the stated full-fibre architecture, not for arbitrary monotone operators.
The quotient construction changes the Banach carrier and preserves the
full-dual certificate, not the underlying mechanism. The parameter and
partner variations describe consequences of that mechanism. None of
these statements identifies independent discoveries merely by counting
the resulting counterexamples.

\else
\section{The mechanism of the counterexample}
\label{sec:structure}
The construction uses the coordinates $r(a)$ and $ma$ to express the
symmetric part of the graph pairing:
\[
\pair{La-Lb}{a-b}=|r(a)-r(b)|^2-\|ma-mb\|_2^2.
\]
A Lipschitz curve in these coordinates makes the graph monotone.
Maximality requires the additional information supplied by the complete
affine fibres. Testing every kernel direction forces an arbitrary polar
query to have the scalar form in Eq.~\eqref{eq:abs-polar-square}.
At an exterior parameter this form returns an actual graph point.
Over the middle interval, the two endpoint comparisons force the
completed profile $(\tau,\eta)$.

Theorem~\ref{thm:endpoint-criterion} gives the exact distinction:
the completed profile contributes a polar point precisely when
some $p\in X^*$ satisfies both $r(p)=\tau$ and $Mp=(\tau,\eta)$.
For the constructed operator, $mp\in\ell^1$ and
$\eta=(1/(4n))_{n\ge1}\notin\ell^1$. Thus the comparison geometry
has a middle segment with no continuous-dual labels realizing it.
The proof takes limits in the comparison space; it never enlarges
the dual space of $c_0$.

This failure of realization can also be seen along actual graph rows.
Set $t_n^\pm=\pm(1+1/n)$, $a_n^\pm=a^{t_n^\pm}$, and
$x_n^\pm=La_n^\pm$. For $t\in J$ and $\epsilon=\sgn t$,
direct substitution in Eqs.~\eqref{eq:D4}--\eqref{eq:D9} gives
\[
(La^t)_{0,1}=-2t-2\epsilon,\qquad
(La^t)_{0,2}=-3t-2\epsilon,\qquad
(La^t)_{0,3}=-t-\epsilon,
\]
with the remaining block-zero coordinates zero. In block $b\ge1$,
the only nonzero coordinate is $-\omega_{|t|-1}(b)$ at index $1$.
Consequently, $x_n^\pm$ converge in the supremum norm to
$x_\pm\in c_0(I)$ with block-zero coordinates $\mp(4,5,2)$ and
positive-block first coordinates $-1/(4b)$. Indeed, convergence of
the tails in $\ell^2$ implies convergence in the supremum norm.
The limits satisfy $\pair{x_\pm}{g}=\pm1$ and lie outside $\dom A$.
On the other hand,
\[
\|a_n^\pm\|_1=2(1+1/n)+1+\|\omega_{1/n}\|_1\longrightarrow\infty.
\]
For each fixed $N$, the sum of the first $N$ tail coordinates tends to
$\sum_{b=1}^N1/(4b)$; these partial sums are unbounded. This proves
the divergence without asserting convergence of the labels in the dual.

The sum obstruction uses a different consequence of the pairing identity.
Every row of the first graph satisfies
\[
\pair{La}{a}=r(a)^2-1-W(|r(a)|-1)>-\sigma,\qquad
\|La\|_\infty>1/2.
\]
These inequalities prove $\mathcal V(\gra A)\le2\sigma$ directly.
The second factor adds exactly $r(a)^2$, making the pairing positive
on the entire sum graph. Since zero remains outside the domain,
$(0,0)$ is a missing polar point of the sum. The finite radial bound
is thus an output of the construction, not an assumed universal
property or an unproved consequence of the desired counterexample.

\fi

\section{Conclusions}
\label{sec:conclusion}
\ifdefined\expandedversion
We constructed maximally monotone operators on $c_0$ and on standard
$\ell^1$ whose sums with specified positive rank-one full-domain
operators are not maximally monotone. In both spaces, the classical
interior-domain condition holds and $(0,0)$ gives a proper monotone
extension of the sum. The maximality proof combines an exact triangular
pairing identity, complete affine fibres, and finite-coordinate endpoint
limits that force a nonsummable dual label in the only remaining case.
The quotient argument preserves the full dual and transfers the same
sum witness to $\ell^1$. The construction also gives a domain at
positive distance from zero whose convex hull contains zero, together
with the explicit Fitzpatrick split values computed below.
For the same first factors, every positive multiple of the rank-one
operator yields a nonmaximal sum. At small strengths, the strict polar
witness moves away from the origin; the first factor is unchanged.
The countable seed and the extension argument also produce first factors
provably different from the determined ones, with full-domain norm
subdifferentials as partners. Their selection is nonconstructive, but
their premises, global maximality, radial bounds and sum witnesses are
all proved above. Both constructions retain the same distinction between
the completed comparison profile and the available continuous-dual labels.
\else
We constructed a maximally monotone operator on $c_0$ whose sum with a
full-domain positive rank-one operator is not maximally monotone under
the classical interior-domain condition. The proof combines an exact
triangular pairing identity, complete affine fibres and two endpoint
comparisons that would force an unavailable continuous-dual label.
The same identity proves the finite radial bound and verifies the
missing polar point of the sum directly. The abstract criterion
identifies exactly when the endpoint comparison adds a graph point,
separating the Hilbert-space comparison from realization in the
original dual pair.
\fi

\ifdefined\expandedversion
\appendix
\section{Values of the partial infimal convolution at the origin}
\label{sec:split}
\subsection{The $c_0$ split}
\label{sec:C7}
For $\alpha\in\R$, define $H(\alpha)=\Phi_A(0,\alpha g)$.
Actual realization of each $t$, Eq.~\eqref{eq:C2}, and
Eq.~\eqref{eq:C9} yield
\begin{equation}
\begin{aligned}
H(\alpha)&=\sup_{|t|>1}[1+W(|t|-1)-t^2+\alpha t]
 \le1+\sigma+\alpha^2/4,\\
H(\alpha)&\ge\sigma+|\alpha|,\qquad H(0)=\sigma.
\end{aligned}
\tag{C12}\label{eq:C12}
\end{equation}
The lower bound uses both actual endpoint limits. For $p\in Y^*$,
$\Phi_P(0,p)=\sup_x[\pair{x}{p}-\pair{x}{g}^2]$.
If $p\notin\R g$, it does not annihilate $\ker g$: a coordinate outside
the first two, or the sum of those two coordinates, supplies a
finite-support witness. Scaling that witness makes the supremum infinite.
If $p=\beta g$, all scalar values $\pair{x}{g}$ are attainable, giving
$\Phi_P(0,p)=\beta^2/4$.

No $\Phi_A$ value is minus infinity, since its defining graph is nonempty.
Off-span splits therefore contribute $+\infty$. The all-dual infimum is
\[
\Delta_{A,P}(0,0)=\inf_{\beta\in\R}[H(-\beta)+\beta^2/4]=\sigma,
\]
uniquely at $\beta=0$. This refinement is not needed for
Theorem~\ref{thm:C}.

\subsection{The standard $\ell^1$ split}
\label{sec:L6}
For $\alpha\in\R$, Eq.~\eqref{eq:L2}, Eq.~\eqref{eq:L6}, and
$\pair{u}{f}=t$ show $\Phi_T(0,\alpha f)=H(\alpha)$, where $H$ is
Eq.~\eqref{eq:C12}. This uses all rows: their expression depends only on
$t$, and each original $a$ admits a lift.

For $w^*\in\ell^\infty$,
$\Phi_B(0,w^*)=\sup_u[\pair{u}{w^*}-\pair{u}{f}^2]$.
If $w^*$ is not a scalar multiple of $f$, some $u\in\ker f$ satisfies
$\pair{u}{w^*}\ne0$. Otherwise, choose $u_0$ with $\pair{u_0}{f}=1$ and
decompose any $u$ as $(u-\pair{u}{f}u_0)+\pair{u}{f}u_0$, forcing
$w^*=\pair{u_0}{w^*}f$, a contradiction.
Scaling that kernel witness gives $\Phi_B(0,w^*)=+\infty$.
For $w^*=\beta f$, all real values of $\pair{u}{f}$ are attainable and
the value is $\beta^2/4$. Therefore the infimum over all $\ell^\infty$
splits is
\[
\Delta_{T,B}(0,0)=\inf_{\beta\in\R}[H(-\beta)+\beta^2/4]=\sigma,
\]
uniquely at $\beta=0$.

\ifdefined\expandedversion
\section{Inputs and the radial conversion}
We use the notation of the construction above. To state the consequences on both
spaces together, let
\[
(\mathcal X,M,k,P_k)=
\begin{cases}
(c_0(I),A,g,P),&I=\{0,1,\ldots\}\times\N,\\
(\ell^1(\N),T,f,B),&f=Q^*g,
\end{cases}
\qquad P_kx=\pair{x}{k}k.
\]
The duals are respectively the full $\ell^1(I)$ and $\ell^\infty(\N)$.
We retain the gap and split conventions
\[
\Phi_M(z,p)=\sup_{(x,a)\in\gra M}\pair{z-x}{a-p},\qquad
\Delta_{M,S}(z,p)=\inf_{b\in\mathcal X^*}
       [\Phi_M(z,p-b)+\Phi_S(z,b)].
\]
Theorems~\ref{thm:C} and~\ref{thm:L} and
Eqs.~\eqref{eq:C2}, \eqref{eq:C9}, and~\eqref{eq:L6} prove
that $M$ is maximally monotone and that every $(x,a)\in\gra M$ satisfies
\begin{equation}
\begin{gathered}
t=\pair{x}{k},\qquad |t|>1,\qquad
\|k\|=2,\qquad \|x\|>\tfrac12,\\
\pair{x}{a}=t^2-1-W(|t|-1),\qquad
0\le W(s)\le\sigma=\frac{\pi^2}{96}<\frac18.
\end{gathered}
\label{eq:cor-input}
\end{equation}
Here $W(s)=\sum_{n\ge1}(\max\{1-sn^{1/4},0\}/(4n))^2$ for $s>0$.
Every $|t|>1$ is realized by an actual graph point: on $c_0(I)$ use the
label $a^t$ in Eq. (D9); on $\ell^1$ use an actual lift through the
surjection $Q$. No bounded linear right inverse is used. The preceding construction
also proves maximality of $P_k$ and the exact split value
\begin{equation}
\Delta_{M,P_k}(0,0)=\sigma.
\label{eq:cor-split-input}
\end{equation}
Only the final finite-split assertion in Corollary~\ref{cor:universal}
uses Eq.~\eqref{eq:cor-split-input}.

For any maximal $M$ on a real Banach space, write
\begin{equation}
L_M(z,p)=\max\left\{0,
\sup_{\substack{(x,a)\in\gra M\\x\ne z}}
\frac{\pair{p-a}{x-z}}{\|x-z\|}\right\}.
\label{eq:cor-Verona}
\end{equation}
We use the nonnegative convention of \cite{verona2008}. The displayed
definition in \cite[Section~2]{verona2012} omits the outer maximum with
zero. At a missing zero primal, the notation in
the preceding construction satisfies exactly
\begin{equation}
\mathcal V(\gra M)=
\sup_{(x,a)\in\gra M}\frac{\max\{0,-\pair{x}{a}\}}{\|x\|}
=L_M(0,0).
\label{eq:cor-V}
\end{equation}
The definition is not a finiteness theorem. For the two concrete
operators, Eq.~\eqref{eq:cor-input} gives, on every graph row,
\[
\pair{x}{a}>-\sigma,\qquad
\frac{\max\{0,-\pair{x}{a}\}}{\|x\|}
\le\frac{\sigma}{\|x\|}<2\sigma.
\]
Taking the supremum proves $\mathcal V(\gra M)\le2\sigma<1/4$.
This bound comes from the construction, not from a universal finiteness
claim or an assumed failure of the sum conjecture.

\begin{lemma}[Radial conversion]\label{lem:radial-conversion}
Let $M:\mathcal X\rightrightarrows\mathcal X^*$ be maximally monotone
on a real Banach space, with $0\notin\dom M$. For $\lambda>0$, put
$J_\lambda=\partial(\lambda\|\cdot\|)$. Then
\begin{equation}
(0,0)\in\gra(M+J_\lambda)^\mu\setminus\gra(M+J_\lambda)
\quad\Longleftrightarrow\quad
\mathcal V(\gra M)\le\lambda.
\label{eq:cor-exact-radial}
\end{equation}
Whenever this holds, $(M,J_\lambda)$ satisfies the original intersection
condition and its sum is not maximally monotone.
\end{lemma}
\begin{proof}
The norm subdifferential has the formula
\[
J_\lambda x=\{j\in\mathcal X^*: \|j\|\le\lambda,
                         \pair{x}{j}=\lambda\|x\|\}.
\]
Hahn--Banach makes this set nonempty for every $x$; at zero it is the
closed dual ball of radius $\lambda$. It is maximally monotone by the
single-subdifferential theorem \cite[Theorem~A]{rockafellar1970subdiff}.
For every $(x,a)\in\gra M$ and every $j\in J_\lambda x$,
\[
\pair{x}{a+j}=\pair{x}{a}+\lambda\|x\|.
\]
Thus relatedness to $(0,0)$ is equivalent to the right-hand side being
nonnegative on every row, which is exactly
$\mathcal V(\gra M)\le\lambda$. The target is outside the sum graph
because $M(0)=\varnothing$. The sum is monotone, so adjoining that target
gives a proper monotone extension. Finally, $\dom J_\lambda=\mathcal X$
and a maximal monotone graph is nonempty, giving the intersection condition.
\end{proof}
The mechanism in this lemma is the missing-primal specialization of the
radial argument in \cite[Theorem~1.2, proof]{verona2008}; the displayed
proof records the full graph quantifiers and the counterexample checks.

\section{Corollaries and their witnesses}
\begin{corollary}[Norm-subdifferential partners]\label{cor:radial}
For each of the two concrete operators $M$, and every $\lambda\ge2\sigma$,
the pair $(M,J_\lambda)$ consists of maximally monotone operators,
$\dom J_\lambda=\mathcal X$, and $M+J_\lambda$ is not maximally monotone.
The target is $(0,0)$. In particular, $\lambda=1/4$ is permitted.
\end{corollary}
\begin{proof}
Use Eq.~\eqref{eq:cor-V} and Lemma~\ref{lem:radial-conversion}.
The finite bound applies to the whole graph of each operator, including
all kernel translates in the $\ell^1$ construction.
\end{proof}

\begin{corollary}[Failure of the regularity conditions]\label{cor:regularity}
For each of the two concrete operators, all five of the following
conditions fail:
\begin{enumerate}
\item $L_M(z,p)=\dist(p,Mz)$ for every $(z,p)\in\mathcal X\times\mathcal X^*$;
\item $M+S$ is maximally monotone for every bounded-range maximally
monotone $S:\mathcal X\rightrightarrows\mathcal X^*$;
\item $M+\partial(\lambda\|\cdot-z\|)$ is maximally monotone for every
$z\in\mathcal X$ and every $\lambda>0$;
\item $L_M(z,0)<\infty$ implies $z\in\dom M$ for every $z\in\mathcal X$;
\item $L_M(z,0)=\dist(0,Mz)$ for every $z\in\mathcal X$.
\end{enumerate}
These are the nonnegative-convention formulation of the regularity
conditions in \cite[Theorem~2.3]{verona2012}. Their displayed definition
uses a bare supremum; the two conventions agree at our missing-zero
witness, since maximality forces a graph row with negative self-pairing.
We prove each failure directly below. In particular, both operators are
nonregular in the convention of \cite{verona2008}.
\end{corollary}
\begin{proof}
We use $\dist(p,\varnothing)=+\infty$. The first, fourth, and fifth
conditions fail at $z=0$, $p=0$, since
\[
L_M(0,0)\le2\sigma<\infty=\dist(0,M0),\qquad 0\notin\dom M.
\]
For the second and third take $S=J_{1/4}$ and $z=0$ in
Corollary~\ref{cor:radial}. Its range is contained in the closed dual
ball of radius $1/4$. This proves each failure directly. Bounded range
is not being inferred from boundedness of the rank-one linear map $P_k$.
\end{proof}

\begin{corollary}[An explicit finite-gap FPV failure]\label{cor:fpv}
For each concrete $M$, set
\begin{equation}
U=\{x\in\mathcal X:\pair{x}{k}>-1\},\qquad
d=0,\qquad p=-\sigma k.
\label{eq:cor-FPV-witness}
\end{equation}
Then $U$ is open and convex, $0\in U$, $U\cap\dom M\ne\varnothing$,
and $(0,p)\notin\gra M$, whereas
\begin{equation}
\pair{-x}{p-a}>0
\quad\text{for every }(x,a)\in\gra M\text{ with }x\in U,
\qquad
\Phi_M(0,p)=2\sigma>0.
\label{eq:cor-FPV-checks}
\end{equation}
Consequently, each of the same operators fails type (FPV) and its
restriction to points at which the Fitzpatrick gap is finite.
\end{corollary}
\begin{proof}
Continuity of $k$ proves openness and convexity. An actual row with
$t=2$ proves $U\cap\dom M\ne\varnothing$. For any graph row in $U$,
$t>-1$ and $|t|>1$ force $t>1$. Therefore
\begin{align*}
\pair{-x}{p-a}
&=\pair{x}{a-p}=t^2-1-W(t-1)+\sigma t\\
&\ge t^2-1+\sigma(t-1)>0.
\end{align*}
On the entire graph, including the negative branch, put $r=|t|>1$.
Then
\begin{align*}
\pair{-x}{a-p}
&=1+W(r-1)-t^2-\sigma t\\
&\le1+\sigma-r^2+\sigma r\\
&=2\sigma-(r-1)(r+1-\sigma)<2\sigma.
\end{align*}
This proves the upper bound on the supremum. Along actual parameters
$t=-1-s$, $s\downarrow0$, we have $W(s)\to\sigma$, and the displayed
expression tends to $2\sigma>0$. Dominated convergence applies since
each squared coordinate is bounded by $1/(16n^2)$.
Finally, $M(0)=\varnothing$ gives graph nonmembership. Type (FPV)
requires membership whenever an open convex set meets the domain and
the target is locally monotonically related there. The finite-gap
restriction imposes in addition $\Phi_M(d,p)<\infty$. All these premises
have just been checked, so both properties fail.
\end{proof}
This witness meets the actual domain. It is not a vacuous test in a
domain-free ball, and no replacement of the first operator by a product
operator has been made.

\begin{corollary}[A bounded-domain normal-cone partner]\label{cor:normal}
For either concrete $M$, choose an actual primal $x_2\in\dom M$ with
$\pair{x_2}{k}=2$. Fix $\rho>\|x_2\|$ and define
\begin{equation}
C_\rho=\{x\in\mathcal X:\|x\|\le\rho,\ \pair{x}{k}\ge-1/2\}.
\label{eq:cor-cut}
\end{equation}
Then $M$ and $N_{C_\rho}$ are maximally monotone,
$\dom M\cap\operatorname{int}\dom N_{C_\rho}\ne\varnothing$, and
their common domain is bounded. Their sum is not maximally monotone;
$(0,-\sigma k)$ belongs to the monotone polar of the sum graph
and lies outside that graph.
In both spaces one can take $\|x_2\|=8$ and $\rho=9$.
\end{corollary}
\begin{proof}
The actual choices from the core proofs are
\[
x_2=-6e_{0,1}-8e_{0,2}-3e_{0,3}\quad\text{on }c_0(I),
\qquad x_2=8e_{n_0}\quad\text{on }\ell^1,
\]
where $q_{n_0}=-(3/4)e_{0,1}-e_{0,2}-(3/8)e_{0,3}$. Both have norm
$8$, and their pairing with $k$ is $2$. The set $C_\rho$ is closed,
convex, bounded, and contains both $0$ and
$x_2$ in its interior. The normal cone is
\[
N_{C_\rho}(x)=\{n\in\mathcal X^*:
\pair{y-x}{n}\le0\ \text{for all }y\in C_\rho\}
\quad(x\in C_\rho),
\]
and is empty off $C_\rho$. It is the subdifferential of the proper
lower semicontinuous convex indicator, hence is maximally monotone
\cite[Theorem~A]{rockafellar1970subdiff}. Since zero is a normal at
every point of $C_\rho$, its domain is exactly $C_\rho$.
Thus $x_2$ verifies the intersection condition, and the common domain
is contained in the ball of radius $\rho$.

Let $p=-\sigma k$. For every $(x,a)\in\gra M$ with $x\in C_\rho$,
$t\ge-1/2$ and $|t|>1$ imply $t>1$. For every $n\in N_{C_\rho}(x)$,
testing the normal inequality at $y=0$ gives $\pair{x}{n}\ge0$.
The calculation in Corollary~\ref{cor:fpv} now gives
\[
\pair{-x}{p-a-n}=\pair{x}{a-p}+\pair{x}{n}>0.
\]
It holds for every sum value. The target is not in the sum graph because
$0\notin\dom M$, so it gives a proper monotone extension.
\end{proof}
This supplies a concrete bounded-common-domain failure, the restricted
sum assertion shown equivalent to the unrestricted one in
\cite[Theorem~2]{botbuenosimons2016}. No maximality of an intermediate
sum is assumed, and boundedness is not replaced by compactness.

\begin{corollary}[Arbitrarily small positive slope bounds]\label{cor:flat}
For either space and each $\theta>0$, there is a maximally monotone
$M_\theta$ with $0\notin\dom M_\theta$ such that
\[
0<\mathcal V(\gra M_\theta)\le\theta.
\]
For each prescribed $\lambda>0$, one may choose such an operator so
that $M_\theta+J_\lambda$ is a nonmaximal sum satisfying the original
intersection condition. The zero endpoint is impossible for a maximal
operator with missing zero primal.
\end{corollary}
\begin{proof}
For $s>0$, define $M_s(x)=sM(x)$. Dividing a polar query's dual
coordinate by $s$ proves that $M_s$ is maximally monotone; its domain
is unchanged. Direct homogeneity gives
$\mathcal V(\gra M_s)=s\mathcal V(\gra M)$.
Taking $s=\theta/(2\sigma)$ proves the upper bound. A zero value would
make $\pair{x}{a}\ge0$ on the entire graph, so $(0,0)$ would be in
its polar; maximality would put it in the graph, contradicting the
missing zero primal. For a prescribed $\lambda$, take
$\theta=\lambda/2$ and apply Lemma~\ref{lem:radial-conversion}.
\end{proof}
Here the quantifier is every positive bound with a possibly different
scaled operator. It does not assert that one fixed operator has slope
below every positive number. In particular, one may take $\theta=1/2$.

\begin{corollary}[Universal assertions and their negations]\label{cor:universal}
Conditional on the two core theorems, the following universal assertions
are false, already on each of the two displayed spaces:
\begin{enumerate}
\item Every qualified pair of maximally monotone operators has a maximally
monotone sum, including the restriction to bounded common domain.
\item Every maximally monotone operator is regular, or satisfies any of
the five equivalent conditions in Corollary~\ref{cor:regularity}.
\item Every maximally monotone operator has the local graph-membership
property defining type (FPV), even if that implication is required only
at points with finite $\Phi_M(d,p)$, with or without the requirement
$U\cap\dom M\ne\varnothing$ in its premises.
\item Every maximally monotone operator omitting zero has
$\mathcal V(\gra M)=+\infty$.
\item For any prescribed $\lambda>0$, all maximally monotone operators
have a maximally monotone sum with $J_\lambda$.
\item There is no finite-$\Delta$ counterexample under the original
intersection condition.
\end{enumerate}
The corresponding assertions over all separable real Banach spaces,
or over all real Banach spaces, are therefore false as well.
\end{corollary}
\begin{proof}
The witnesses are, respectively, Corollary~\ref{cor:normal},
Corollary~\ref{cor:regularity}, Corollary~\ref{cor:fpv},
Eq.~\eqref{eq:cor-V} with the finite bound derived after it, and
Corollary~\ref{cor:flat}. For the last assertion, use the original
rank-one pairs and Eq.~\eqref{eq:cor-split-input}. Both spaces are
separable, so the same witnesses disprove the larger-carrier assertions.
\end{proof}

\noindent\textit{Scope of equivalence and attribution.}
Negating $\forall M\,\mathcal P(M)$ produces $\exists M\,\neg\mathcal P(M)$,
not $\forall M\,\neg\mathcal P(M)$. A valid equivalence between universal
assertions remains valid when those assertions are false. The five
fixed-operator equivalences cited above are classical. Each failure has
the explicit witness given above. These corollaries use no additional
equivalence between assertions on different operators or spaces.

\section{Consequences for quotient norms and graph representations}
The following consequences use the actual maps of the construction. They
identify the information retained by the dual quotient and the information
lost by extending a curve in the Hilbert comparison space. They are separate deductions,
not additional assumptions in the core maximality proof.

\begin{corollary}[Constant pairing and an unbounded norm in $\ell^\infty/c_0$]\label{cor:corona}
For the standard $\ell^1$ operator $T$ in the construction above, there are
actual rows $(u_N,p_N)\in\gra T$ and a fixed $d\in\ell^1$ such that
\begin{equation}
\|u_N\|_1=8,\qquad \pair{u_N}{p_N}=3,\qquad
\dist_{\infty}(p_N,c_0)=\|p_N\|_\infty=5+2N,
\label{eq:cor-constant-work}
\end{equation}
and
\begin{equation}
\|d\|_1=\frac{\pi^2}{6},\qquad \pair{d}{f}=0,\qquad
\pair{d}{p_N}=-2\sum_{r\ge1}\frac{\min\{N,r\}}{r^2}\longrightarrow-\infty.
\label{eq:cor-fixed-detector}
\end{equation}
The fibre of $(\gra T)^\mu$ above $d$ is empty. In particular, no finite
constants $K,\eta\ge0$ satisfy
\begin{equation}
\dist_{\infty}(p,c_0)\le K+\eta\pair{u}{p}
\quad\text{for every }(u,p)\in\gra T.
\label{eq:cor-WCC-failure}
\end{equation}
\end{corollary}
\begin{proof}
The quotient norm is $\|p+c_0\|_{\ell^\infty/c_0}=\dist_\infty(p,c_0)$.
We first compute it for the enumeration defining $Q$.
For every $a\in\ell^1(I)$,
\begin{equation}
\dist_{\infty}(Q^*a,c_0)
=\limsup_j|\pair{q_j}{a}|=\|a\|_1.
\label{eq:cor-quotient-isometry}
\end{equation}
Indeed, the upper bound follows from $\|q_j\|_\infty\le1$. For
$a\ne0$ and $0<\varepsilon<\|a\|_1$, a finitely supported rational
vector in the unit ball gives pairing greater than
$\|a\|_1-\varepsilon$ in absolute value. Rational scalar perturbations
tending to one supply infinitely many distinct such vectors (start
with a strict bound). Every one occurs in the enumeration, so the
same lower bound holds after deleting any finite prefix. Let
$\varepsilon\downarrow0$. The case $a=0$ is immediate. For any bounded
sequence, its distance to $c_0$ equals the limsup of the absolute
coordinates: truncation gives one inequality, and subtracting any
norm-null sequence gives the other. This proves
Eq.~\eqref{eq:cor-quotient-isometry}, including
$\|Q^*a\|_\infty=\|a\|_1$.

Use actual labels
\[
a_N=-2e_{0,1}+3e_{0,3}
       +\sum_{b=1}^N(e_{b,1}-e_{b,2}).
\]
They have $r(a_N)=2$ and $ma_N=(1,0,\ldots)=F(2)$, since
$\omega_1=0$. Thus $a_N\in D$. Directly from the triangular map,
\[
x_N=La_N=-6e_{0,1}-8e_{0,2}-3e_{0,3}
                  +\sum_{b=1}^N(e_{b,1}+e_{b,2}).
\]
Consequently $\|x_N\|_\infty=8$, $\pair{x_N}{a_N}=3$, and
$\|a_N\|_1=5+2N$. Let $i_N$ be the first index with
$q_{i_N}=x_N/8$, and put $u_N=8e_{i_N}$ and $p_N=Q^*a_N$.
These are actual rows because $Qu_N=x_N$. The pairing is preserved,
and Eq.~\eqref{eq:cor-quotient-isometry} gives
Eq.~\eqref{eq:cor-constant-work}.

For $r\ge1$, the finite primal vector
$v_r=\sum_{b=1}^r(e_{b,1}-e_{b,2})$ is rational and has norm one.
Let $j_r$ be its first index in the enumeration, so the $j_r$ are
distinct. Set $d=-\sum_{r\ge1}r^{-2}e_{j_r}$. This converges in
$\ell^1$, has norm $\pi^2/6$, and
\[
f(j_r)=\pair{v_r}{g}=0,\qquad
p_N(j_r)=\pair{v_r}{a_N}=2\min\{N,r\}.
\]
For each fixed $N$ the series defining $\pair{d}{p_N}$ converges
absolutely. Its negative is at least
$2\sum_{r=1}^N1/r$, proving Eq.~\eqref{eq:cor-fixed-detector}.
Every primal in $\dom T$ has $|\pair{u}{f}|>1$, whereas
$\pair{d}{f}=0$. Thus $d\notin\dom T$, and maximality
$(\gra T)^\mu=\gra T$ makes its entire polar fibre empty.
Finally, Eq.~\eqref{eq:cor-constant-work} would turn
Eq.~\eqref{eq:cor-WCC-failure} into $5+2N\le K+3\eta$ for every
$N$, which is impossible.
\end{proof}
The fixed vector $d$ detects divergence of the dual sequence, although
all selected graph points have the same primal norm and pairing.
Since $d\notin\dom T$, no point of $(\gra T)^\mu$ has primal coordinate
$d$. In particular, $\dist_\infty(p,c_0)$ cannot be bounded on $\gra T$
by a function of $\|u\|_1$ and $\pair{u}{p}$ that is finite at $(8,3)$.

\begin{corollary}[A Lipschitz extension with no additional preimages in the dual space]\label{cor:dense-profile}
For the displayed $c_0(I)$ construction, the map $(r,m)$ has range
\begin{equation}
\operatorname{range}(r,m)=\R\times\ell^1(\{0,1,\ldots\})
\subsetneq\R\times\ell^2(\{0,1,\ldots\}),
\label{eq:cor-dense-range}
\end{equation}
which is dense. The profile $F:J\to\ell^2$ has a unique global
$1$-Lipschitz extension, namely
\begin{equation}
\widehat F(t)=
\begin{cases}
F(t),&|t|>1,\\
(t,\omega_0),&|t|\le1.
\end{cases}
\label{eq:cor-profile-extension}
\end{equation}
Its graph strictly extends $\gra F$, and is maximal for the quadratic
positivity inequality $\|y-y'\|_2^2\le|t-t'|^2$ on
$\R\times\ell^2$. Nevertheless, its entire actual pullback satisfies
\begin{equation}
\{(La,a):a\in\ell^1(I),\ ma=\widehat F(r(a))\}
=\gra A.
\label{eq:cor-same-pullback}
\end{equation}
For these maps, a maximal monotone graph can therefore have a nonmaximal
image under the quadratic representation described in the proof, despite
the dense range of $(r,m)$.
\end{corollary}
\begin{proof}
Every $ma$ is summable. Conversely, for any $t\in\R$ and any
$y\in\ell^1(\{0,1,\ldots\})$, the actual label
\[
a=-t e_{0,1}+(t+y_0)e_{0,3}+\sum_{n\ge1}y_n e_{n,1}
\]
satisfies $r(a)=t$ and $ma=y$. This proves
Eq.~\eqref{eq:cor-dense-range}; strictness follows from
$\omega_0\in\ell^2\setminus\ell^1$, and density from finite truncations.

The coordinatewise positive-part estimate gives
\[
\|\omega_s-\omega_q\|_2^2\le c|s-q|^2,\qquad
c=\frac1{16}\sum_{n\ge1}n^{-3/2}<1,
\]
also for $s=0$ or $q=0$. Put $\chi(t)=\max\{-1,\min\{1,t\}\}$ and
$\zeta(t)=\max\{|t|-1,0\}$. Then
$\widehat F(t)=(\chi(t),\omega_{\zeta(t)})$.
The same-branch and opposite-branch cases give
\[
|\chi(t)-\chi(u)|+|\zeta(t)-\zeta(u)|\le|t-u|.
\]
Therefore
\[
\|\widehat F(t)-\widehat F(u)\|_2^2
\le|\chi(t)-\chi(u)|^2+c|\zeta(t)-\zeta(u)|^2
\le|t-u|^2.
\]
Any global $1$-Lipschitz extension must take the values
$(-1,\omega_0)$ and $(1,\omega_0)$ at the two endpoints, by
continuity from the actual branches. Their distance is $2$.
For an intermediate $t$, the bounds $t+1$ and $1-t$ on the two
distances add to $2$. Equality in the Hilbert-space triangle inequality
forces the affine segment value $(t,\omega_0)$. Hence the extension
is unique. Its graph is maximal for the displayed quadratic inequality:
testing any related $(t,y)$ against $(t,\widehat F(t))$ forces
$y=\widehat F(t)$.

No new parameter $|t|\le1$ has an actual label, because
$(t,\omega_0)\notin\ell^1$. Thus the label constraint for
$\widehat F$ is exactly the original constraint for $F$, proving
Eq.~\eqref{eq:cor-same-pullback}.

We make the last correspondence assertion precise. Set $V=\R$,
$H=\ell^2(\{0,1,\ldots\})$, and $\mathcal T v=vh$. The core identities are
\begin{align*}
\pair{La}{b}+\pair{Lb}{a}&=2[r(a)r(b)-(ma,mb)_H],\\
\pair{\mathcal T v}{a}&=v r(a),\qquad
(\ker r\cap\ker m)^\perp\cap c_0(I)=\R h.
\end{align*}
For the last identity, annihilating differences within every positive
block makes each block constant; the $c_0$ condition makes it zero.
In block zero the tail similarly vanishes and the first two entries
are equal. The converse follows from $r(k)=0$ on the kernel.
Thus the bilinear and adjoint identities and the displayed annihilator
identity all hold.
On $x=Lp+2\delta h$ the coordinate map is
\[
J_{\rm coord}(x,p)=(r(p)+\delta,mp,\delta),\qquad
q(t,y,\delta)=t^2-\|y\|_2^2-\delta^2.
\]
For this quadratic form, a nonempty set $C$ is $q$-positive if
$q(c-c')\ge0$ for all $c,c'\in C$. Define
$C^q=\{z:q(z-c)\ge0\text{ for every }c\in C\}$.
Maximality is with respect to inclusion among $q$-positive sets.
The image of the actual maximal graph is
$C^0=\{(t,F(t),0):t\in J\}$. It is not maximally $q$-positive,
since $\{(t,\widehat F(t),0):t\in\R\}$ properly extends it and
is $q$-positive. In fact, the entire quadratic polar is exactly
\[
(C^0)^q=\{(t,\widehat F(t),0):t\in\R\}.
\]
To verify the reverse inclusion, let $(\tau,y,\delta)$ be related
to $C^0$. If $|\tau|>1$, testing at the same parameter forces
$y=F(\tau)$ and $\delta=0$. If $|\tau|\le1$, limits along the
actual branches give
\[
\|y-(1,\omega_0)\|^2+\delta^2\le(1-\tau)^2,\qquad
\|y-(-1,\omega_0)\|^2+\delta^2\le(1+\tau)^2.
\]
Multiplication by $(1+\tau)/2$ and $(1-\tau)/2$, respectively,
and addition cancel the common term $1-\tau^2$, leaving
$\|y-(\tau,\omega_0)\|^2+\delta^2\le0$.
This also proves uniqueness of the maximal reduced extension, even
when the coordinate $\delta$ is allowed. Nevertheless its new points
have no actual preimages under $J_{\rm coord}$.
Consequently, dense joint range cannot replace actual
joint surjectivity in a correspondence sending every actual maximal
extension to a maximal reduced positive set. This is a failure for
the displayed maps, not an assertion excluding every alternative
representation of $A$.
\end{proof}
The Lipschitz extension is unique. Its values at $|t|\le1$ have no
preimages under $(r,m)$, so the graph defined by the displayed
constraint remains unchanged.

\begin{corollary}[A gap between optimization over the domain and its convex hull]\label{cor:two-row}
On each of the two original spaces there is an operator $\widetilde M$
obtained from $M$ by an affine change of variables, a norm-one functional $\varphi$, and a
unit vector $e$ with $\pair{e}{\varphi}=1$ such that
\begin{equation}
0\le \sup_{\substack{z\in\dom\widetilde M\\\|z\|\le1}}
\pair{z}{\varphi}\le\frac38,
\qquad
\inf_{p\in X^*}\left\{
\sup_{z\in\dom\widetilde M}\pair{z}{p}+\|\varphi-p\|
\right\}=1.
\label{eq:cor-PCT-gap}
\end{equation}
There are two graph points whose mean primal coordinate is $3e/4$.
Its value under $\varphi$ is $3/4$, strictly above the left supremum in
Eq.~\eqref{eq:cor-PCT-gap}. Moreover,
\begin{equation}
\dist(3e/4,\dom\widetilde M)\ge\frac38,
\qquad \Phi_{\widetilde M}(3e/4,0)=\sigma.
\label{eq:cor-finite-projection-hole}
\end{equation}
The transformation changes the operator, but not the Banach space or
its norm.
\end{corollary}
\begin{proof}
Let $x_2$ denote the actual primal of norm $8$ used above, with
$\pair{x_2}{k}=2$. Both $x_2$ and $-x_2$ are actual primals.
Set $\varphi=k/2$. On $c_0(I)$ take
$e=e_{0,1}-e_{0,2}$, regarded as a primal vector. On standard
$\ell^1$ take $e=e_j$ at an index where
$q_j=e_{0,1}-e_{0,2}$. In both cases
$\|e\|=\|\varphi\|=1$ and $\pair{e}{\varphi}=1$.
Put $r_0=3/4$ and define
\[
\mathcal S z=-\frac{x_2}{r_0}\pair{z}{\varphi}
                 +z-\pair{z}{\varphi}e,
\qquad
\widetilde M(z)=\mathcal S^*M(x_2+\mathcal S z).
\]
Here $\mathcal S$ is a bounded linear isomorphism, and the dual variable
is transformed by $\mathcal S^*$. Indeed, writing $z=\alpha e+w$ with $w\in\ker\varphi$
gives $\mathcal S z=-\alpha x_2/r_0+w$ and
$\pair{\mathcal S z}{\varphi}=-\alpha/r_0$. For a prescribed
image $v$, the unique inverse is obtained with
$\alpha=-r_0\pair{v}{\varphi}$ and
$w=v+\alpha x_2/r_0\in\ker\varphi$; it is bounded.
The identity $\pair{z-z'}{\mathcal S^*(a-a')}
=\pair{\mathcal S(z-z')}{a-a'}$ preserves every monotonicity bracket.
The resulting bijection on primal--dual pairs therefore preserves
monotonicity and full maximality.

Every $z\in\dom\widetilde M$ satisfies
\[
1<|\pair{x_2+\mathcal S z}{k}|
  =|2-2\pair{z}{\varphi}/r_0|.
\]
Consequently
\begin{equation}
\pair{z}{\varphi}<\frac38
\quad\text{or}\quad
\pair{z}{\varphi}>\frac98.
\label{eq:cor-transformed-strip}
\end{equation}
The second alternative is impossible when $\|z\|\le1$.
The primals $0$ and $2r_0e=3e/2$ are actual, since their images
are $x_2$ and $-x_2$. This proves the singleton bounds in
Eq.~\eqref{eq:cor-PCT-gap} and provides strict access inside the
unit ball at $0$.

Writing $D'=\dom\widetilde M$, the actual endpoints imply
$e\in\operatorname{conv}D'$. For every $p\in X^*$,
\[
\sup_{z\in D'}\pair{z}{p}+\|\varphi-p\|
\ge\pair{e}{p}+\pair{e}{\varphi-p}=1.
\]
The choice $p=0$ attains $1$, proving the second equality in
Eq.~\eqref{eq:cor-PCT-gap}. Choose graph points over $0$ and $3e/2$
and assign each weight $1/2$. Their mean primal coordinate is $3e/4$,
which lies strictly inside the unit ball. Its value under $\varphi$
is $3/4$, exceeding the supremum on the left of
Eq.~\eqref{eq:cor-PCT-gap}.

For $d'=r_0e$ one has $x_2+\mathcal S d'=0$.
Eq.~\eqref{eq:cor-transformed-strip} and $\|\varphi\|=1$
give the distance lower bound. The definition of the Fitzpatrick
function and the full bijection of graph rows give
\[
\Phi_{\widetilde M}(d',0)
=\sup_{(x,a)\in\gra M}\pair{-x}{a}
=\Phi_M(0,0)=\sigma.
\]
This proves Eq.~\eqref{eq:cor-finite-projection-hole}.
\end{proof}
The two quantities in Eq.~\eqref{eq:cor-PCT-gap} are computed directly.
Equation~\eqref{eq:cor-finite-projection-hole} also gives a point
$(3e/4,0)$ where the Fitzpatrick function is finite although $3e/4$
lies outside the norm closure of the domain.

\begin{corollary}[Convex combinations of graph points and incompatible polar points]
\label{cor:incompatible}
For each original rank-one pair $(M,P_k)$, put
$G=\gra(M+P_k)$. We construct two points of $\gra M$ and one point
of $\gra P_k$ with the same mean primal coordinate. Their pairings produce
a point $q'$ and an explicit $\eta>0$ such that
\begin{equation}
q_0=(0,0)\in G^\mu,\qquad
[q',g]\ge\eta\quad(g\in G),\qquad
[q_0,q']=-\eta.
\label{eq:cor-incompatible}
\end{equation}
Here $[(x,a),(y,b)]=\pair{x-y}{a-b}$.
Thus $G^\mu$ is not monotone, and $G$ has at least two distinct,
incomparable maximally monotone extensions.
\end{corollary}
\begin{proof}
Set $t=65/64$, $s=1/64$, and
\[
v=\sum_{n\ge1}\omega_s(n)e_{n,1},\qquad
W=\|\omega_s\|_2^2,\qquad E=t^2-1=\frac{129}{4096}.
\]
The vector $v$ is finitely supported. The single first tail coordinate
already gives
\[
W\ge\left(\frac{63}{256}\right)^2
  =\frac{3969}{65536}>\frac{2064}{65536}=E.
\]
Take the two actual labels
\[
a^+=-t e_{0,1}+(t+1)e_{0,3}+v,
\qquad
a^-=t e_{0,1}-(t+1)e_{0,3}+v.
\]
They realize parameters $t$ and $-t$. On $c_0(I)$ put
$x^\pm=La^\pm$; their equal-weight mean is
$w=Lv=-\sum_n\omega_s(n)e_{n,1}$, and their dual mean is $v$.
One has $\pair{w}{g}=0$, so $(w,0)$ is an actual $P$-row.
Moreover,
\[
\pair{x^\pm}{a^\pm}=E-W,\qquad
\pair{w}{v}=-W.
\]
The correction term for the mean of the two points of $\gra A$ is
\[
\frac14\pair{x^+-x^-}{a^+-a^-}
=\frac12\pair{x^+}{a^+}+\frac12\pair{x^-}{a^-}
  -\pair{w}{v}=E.
\]
The corresponding correction for the single point of $\gra P$ is zero,
and $\pair{-w}{v}=W>E$.

On standard $\ell^1$, choose actual lifts $u^\pm$ with
$Qu^\pm=x^\pm$, and use rows $(u^\pm,Q^*a^\pm)$.
Their mean $\bar u=(u^++u^-)/2$ has
$\pair{\bar u}{f}=0$, so $(\bar u,0)$ is an actual $B$-row.
All the preceding pairing equalities are preserved by
$\pair{u}{Q^*a}=\pair{Qu}{a}$. No particular choice of lift,
linear right inverse, or change of norm is needed.

Write $r=(w,v)$ in the first space and
$r=(\bar u,Q^*v)$ in the second. Averaging monotonicity with
the specified graph points proves
\[
[r,g]\ge-E\quad(g\in G),\qquad [q_0,r]=-W.
\]
Indeed, for a sum row $(x,a+b)$, the averaged $M$ brackets
give $\pair{r_X-x}{r_{X^*}-a}\ge-E$, and the single
partner bracket gives $\pair{r_X-x}{-b}\ge0$.
The core relatedness $q_0\in G^\mu$ can also be read directly
from $\pair{x}{a+P_kx}=2|\pair{x}{k}|^2-1-W(|\pair{x}{k}|-1)>0$.
For
\[
\theta=\frac{W-E}{2W},\qquad q'=\theta r,
\qquad \eta=\frac{(W-E)^2}{4W}>0,
\]
the three-point identity yields
\begin{align*}
[q',g]
&=(1-\theta)[q_0,g]+\theta[r,g]
                 -\theta(1-\theta)[q_0,r]\\
&\ge\theta\bigl((1-\theta)W-E\bigr)=\eta.
\end{align*}
Also $[q_0,q']=\theta^2[r,q_0]=-\eta$.
Thus both points belong to $G^\mu\setminus G$ and no monotone
extension of $G$ can contain both. Each monotone set
$G\cup\{q_0\}$ and $G\cup\{q'\}$ has a maximally monotone
extension: a chain of monotone supersets has a monotone union,
so Zorn's lemma applies. The two extensions differ, and neither
can properly contain the other by maximality.
\end{proof}
The graph points and their pairings are given above; the maximal
extensions in the final sentence are existence conclusions, not
two additional explicit operator formulas. The averaged pair $r$
is not asserted to be a factor row. This distinction is essential:
the finite averaging identity produces the new polar point through
the proved bracket calculation, not by identifying a barycentre
with an actual graph row.

\fi
\fi

\section*{Use of AI}
The project was carried out under human direction, including decisions on
the research objective, source materials, presentation, and review requirements.
AI systems were used substantially for mathematical construction, proof
development and reconstruction, error detection and correction, Lean
formalization, and preparation of the exposition.

\begin{thebibliography}{99}
\raggedright
\bibitem{simons2008}
S. Simons,
\emph{From Hahn--Banach to Monotonicity}, 2nd ed.,
Lecture Notes in Mathematics, vol.~1693, Springer, 2008.
doi:10.1007/978-1-4020-6919-2.

\bibitem{fitzpatrick1988}
S. Fitzpatrick,
Representing monotone operators by convex functions,
in: Workshop/Miniconference on Functional Analysis and Optimization,
Proceedings of the Centre for Mathematical Analysis, Australian National
University, vol.~20, 1988, pp.~59--65.

\bibitem{burachiksvaiter2002}
R.~S. Burachik and B.~F. Svaiter,
Maximal monotone operators, convex functions and a special family of enlargements,
\emph{Set-Valued Analysis} \textbf{10} (2002), 297--316.
doi:10.1023/A:1020639314056.

\bibitem{penot2004}
J.-P. Penot,
The relevance of convex analysis for the study of monotonicity,
\emph{Nonlinear Analysis: Theory, Methods \& Applications}
\textbf{58} (2004), 855--871.
doi:10.1016/j.na.2004.05.018.

\bibitem{martinezsvaiter2005}
J.-E. Mart\'{i}nez-Legaz and B.~F. Svaiter,
Monotone operators representable by l.s.c. convex functions,
\emph{Set-Valued Analysis} \textbf{13} (2005), 21--46.
doi:10.1007/s11228-004-4170-4.

\bibitem{buenomartinezsvaiter2014}
O. Bueno, J.-E. Mart\'{i}nez-Legaz, and B.~F. Svaiter,
On the monotone polar and representable closures of monotone operators,
\emph{Journal of Convex Analysis} \textbf{21} (2014), 495--505.

\bibitem{rockafellar1970}
R.~T. Rockafellar,
On the maximality of sums of nonlinear monotone operators,
\emph{Transactions of the American Mathematical Society}
\textbf{149} (1970), 75--88.

\bibitem{yao2010}
L. Yao,
The sum of a maximal monotone operator of type (FPV) and a maximal
monotone operator with full domain is maximal monotone,
\emph{Nonlinear Analysis: Theory, Methods \& Applications}
\textbf{74} (2011), 6144--6152.
doi:10.1016/j.na.2011.05.093.

\bibitem{borweinyao2012}
J.~M. Borwein and L. Yao,
Maximality of the sum of a maximally monotone linear relation and a
maximally monotone operator,
\emph{Set-Valued and Variational Analysis} \textbf{21} (2013), 603--616.
doi:10.1007/s11228-013-0259-y.

\bibitem{voisei2010}
M.~D. Voisei,
A sum theorem for (FPV) operators and normal cones,
\emph{Journal of Mathematical Analysis and Applications}
\textbf{371} (2010), 661--664.
doi:10.1016/j.jmaa.2010.05.064.

\bibitem{marquessvaiter2012}
M. Marques Alves and B.~F. Svaiter,
A new qualification condition for the maximality of the sum of maximal
monotone operators in general Banach spaces,
\emph{Journal of Convex Analysis} \textbf{19} (2012), 575--589.

\bibitem{gossez1971}
J.-P. Gossez,
Op\'{e}rateurs monotones non lin\'{e}aires dans les espaces de Banach non r\'{e}flexifs,
\emph{Journal of Mathematical Analysis and Applications}
\textbf{34} (1971), 371--395.
doi:10.1016/0022-247X(71)90119-3.

\bibitem{marquessvaiter2010}
M. Marques Alves and B.~F. Svaiter,
On Gossez type (D) maximal monotone operators,
\emph{Journal of Convex Analysis} \textbf{17} (2010), 1077--1088.

\bibitem{bauschkefp2012}
H.~H. Bauschke, J.~M. Borwein, X. Wang, and L. Yao,
Every maximally monotone operator of Fitzpatrick--Phelps type is actually
of dense type,
\emph{Optimization Letters} \textbf{6} (2012), 1875--1881.
doi:10.1007/s11590-011-0383-2.

\bibitem{buenosvaiter2011}
O. Bueno and B.~F. Svaiter,
A non-type (D) operator in $c_0$,
\emph{Mathematical Programming} \textbf{139} (2013), 81--88.
doi:10.1007/s10107-013-0661-0.
Lemma numbering refers to the author version, arXiv:1103.2349v1.

\bibitem{bauschkeetal2011}
H.~H. Bauschke, J.~M. Borwein, X. Wang, and L. Yao,
Construction of pathological maximally monotone operators on
non-reflexive Banach spaces,
\emph{Set-Valued and Variational Analysis} \textbf{20} (2012), 387--415.
doi:10.1007/s11228-012-0209-0.
Example numbering refers to the authors' August~6, 2011 version.

\bibitem{eberhardwenczel2024}
A. Eberhard and R. Wenczel,
Representative functions, variational convergence and almost convexity,
\emph{Set-Valued and Variational Analysis} \textbf{32} (2024), article~4.
doi:10.1007/s11228-024-00708-4.

\bibitem{botcsetnek2008}
R.~I. Bo\c{t} and E.~R. Csetnek,
Weak regularity conditions for maximal monotonicity in separable
Asplund spaces,
Chemnitz Faculty of Mathematics, Preprint~3, April~1, 2008.

\bibitem{verona2008}
A. Verona and M. E. Verona,
Regular maximal monotone multifunctions and enlargements,
\emph{Journal of Convex Analysis} \textbf{16} (2009), 1003--1009.

\bibitem{rockafellar1970subdiff}
R. T. Rockafellar,
On the maximal monotonicity of subdifferential mappings,
\emph{Pacific Journal of Mathematics} \textbf{33} (1970), 209--216.

\bibitem{botbuenosimons2016}
R. I. Bo\c{t}, O. Bueno and S. Simons,
The Rockafellar conjecture and type (FPV),
\emph{Set-Valued and Variational Analysis} \textbf{24} (2016), 381--385.
doi:10.1007/s11228-016-0384-5.
\ifdefined\expandedversion
\bibitem{verona2012}
M. E. Verona and A. Verona,
On the regularity of maximal monotone operators and related results,
arXiv:1212.1968v3, 2012.
\fi
\end{thebibliography}

\end{document}
